\documentclass[11pt,a4paper]{ctexart}

% ---------- 页面 ----------
\usepackage[a4paper,left=2.4cm,right=2.4cm,top=2.4cm,bottom=2.2cm]{geometry}
\usepackage{fancyhdr}
\pagestyle{fancy}
\fancyhf{}
\fancyhead[L]{\small 数值库数学串讲}
\fancyhead[R]{\small\leftmark}
\fancyfoot[C]{\small\thepage}
\renewcommand{\headrulewidth}{0.4pt}

% ---------- 数学 ----------
\usepackage{amsmath,amssymb,amsthm,mathtools}
\usepackage{array,booktabs,longtable}
\usepackage{enumitem}
\usepackage{xcolor}
\usepackage{titlesec}
\usepackage[hidelinks]{hyperref}

\setlist{itemsep=1pt,topsep=2pt,parsep=0pt}
\titleformat{\section}{\Large\bfseries}{\thesection}{0.6em}{}
\titleformat{\subsection}{\large\bfseries}{\thesubsection}{0.6em}{}

% ---------- 条目环境 ----------
% \begin{know}{编号}{名称}{状态}
\newenvironment{know}[3]
  {\par\medskip\noindent\hangindent=0pt
   \textbf{\textsf{#1}}\ \textbf{#2}\quad{\small\textit{#3}}\par\nopagebreak\smallskip}
  {\par\medskip}

\newcommand{\field}[1]{\textbf{#1}\ }
\newcolumntype{L}[1]{>{\raggedright\arraybackslash}p{#1}}

\newcommand{\eps}{\varepsilon}
\newcommand{\C}{\mathbb{C}}
\newcommand{\R}{\mathbb{R}}
\newcommand{\Z}{\mathbb{Z}}
\newcommand{\HH}{\mathbb{H}}
\newcommand{\abs}[1]{\left\lvert #1\right\rvert}
\newcommand{\norm}[1]{\left\lVert #1\right\rVert}
\newcommand{\inner}[2]{\langle #1,#2\rangle}

\title{\bfseries 数值库数学串讲\\[0.3em]
\large 从公式到判据：一份可查漏的数学清单}
\author{}
\date{}

\begin{document}
\maketitle

\noindent
这份文档只写数学，不写代码。它回答两个问题：
\begin{enumerate}[label=\arabic*.]
  \item 这套数值库里，\textbf{每一个数学对象、每一条公式、每一个阈值}分别是什么，
        成立条件是什么，失败了长什么样；
  \item 这些零散的点之间\textbf{靠什么连成一张网}——每节开头的「串讲」小节就是这张网的边。
\end{enumerate}

\begingroup\sloppy
\noindent
所有事实以源码头文件末尾的「三、公式与符号表」，以及 \texttt{01-numerics.md}
（在 \texttt{docs/standard/} 下，D7/D12/D19/D20）为准；凡是原文没给的，
一律标「原文未给」并归入第 \ref{sec:audit} 节的缺口，不用常识替它补。
\par\endgroup

\tableofcontents
\clearpage

\section{总图：十个分支与一张依赖网}

\subsection{状态图例}

全篇每个条目后面都跟一个状态词，取值为：

\begin{center}
\begin{tabular}{L{2.6cm}L{12.4cm}}
\toprule
状态 & 含义 \\
\midrule
已实现 & 源码头文件标注「已实现」，函数体有真实算法 \\
框架 & 签名与注释已定、函数体全部 \texttt{throw}；注释即规格，公式可信，数值未验证 \\
骨架 & 同上，但源码自述为「骨架」；公式仍是规格 \\
未建 & 只在 \texttt{NOW.md} 里规划，文件不存在 \\
缺位 & 库里既没有位置、也没有规划；必须在第 \ref{sec:audit} 节的缺口里登记 \\
\bottomrule
\end{tabular}
\end{center}

\textbf{框架与骨架的数学含量并不低}：它们是「先写契约后写代码」的产物，公式段往往是全仓最完整的。
但它们的\textbf{判据全部未经验证}——条目里凡是标「未验证」的阈值，都只代表意图，不代表实测。

\subsection{十个分支与条目编号}

条目编号是全局的、不按章的先后。同一编号在全文中唯一，交叉引用都用它。

\begin{center}
\begin{tabular}{L{2.5cm}L{2.9cm}L{8.4cm}}
\toprule
分支 & 编号 & 覆盖对象 \\
\midrule
分析基础 & A1--A6 & 机器精度、累加误差、前向/后向误差、条件数、收敛阶、相消 \\
线性代数 & L1--L10 & 线性映射、矩阵乘、转置与伴随、内积、范数与迹、谱、奇异值、正交相似 \\
数值线性代数 & N1--N13 & 部分主元、Householder、Givens、MGS、Cholesky、Padé、Jacobi、$A^{H}A$ 路线、秩判定 \\
复分析 & C1--C6 & 复代数、Smith 除法、欧拉公式、主值与分支切割、开方、幂 \\
调和分析 & H1--H11 & DFT、归一化、卷积定理、频轴、radix-2、Parseval、级数投影、Gibbs、谱收敛 \\
概率 & P1--P10 & 特征函数、矩、接缝、经验特征函数、Box--Muller、erf/$\Phi$、分位数 \\
微分与自动微分 & D1--D8 & 对偶数、链式法则、雅可比布局、种子 sweep、反向模式、Lipschitz、RK 阶、辛结构 \\
几何与李群 & G1--G10 & 叉积与 hat、Rodrigues、李括号、四元数代数、双覆盖、插值、Shepperd \\
离散化与变分 & F1--F9 & 网格与自由度、P1 刚度、弱形式、一致载荷、Dirichlet、收敛阶、CSR、CG \\
张量代数 & T1--T7 & 形状与步长、偏移、连续性、零拷贝视图、reshape 合法性、广播 \\
微分形式 & M1--M7 & $k$-形式、外积、外微分、拉回、Hodge、Stokes、$n=3$ 字典 \\
\bottomrule
\end{tabular}
\end{center}

\subsection{依赖网（ASCII）}

箭头读作「是……的基础」。方括号里是状态。

\begin{verbatim}
 A1 eps --> A2 8n*eps --> A3 fwd/bwd err --> A4 cond number
   |            |               |                |
   |            v               v                v
   |          N1 pivot     N5 MGS ill-cond   N12 A^H A squared
   v            |                                ^
 C2 Smith       v                                |
            N3 Householder --> N6 thin QR    L7 singular values
                 |                 |              ^
                 v                 v              |
               N4 Givens      N6b least sq   N11 Jacobi spectrum

 L1 row-major --> L2 matmul --> L3 transpose/adjoint
   |                              |
   v                              v
 T1 shape/stride --> T2 offset --> T3 contiguous --> T4 view
                                                        |
                                                        v
                                                   T5 reshape
 D1 dual eps^2=0 --> D2 chain rule --> D3 jacobian layout
                                        |--> D4 seed sweep (col)
                                        \--> D5 reverse mode (row)

 H1 DFT --> H2 1/N --> H3 circular conv --> H4 linear conv (pad)
   |          |                                |
   v          v                                v
 H5 freq axis  H8 Parseval            H10 Gibbs / H11 spectral decay
   |
   v
 P1 char func (+i, opposite sign) --> P3 conv theorem --> P6 empirical
                                          |
                                          v
                                     P4 moments = derivatives

 G1 cross/hat --> G2 Rodrigues --> G4 Lie bracket
      |               |
      v               v
 G6 quat double cover --> G8 slerp --> G9 Shepperd

 D6 Lipschitz --> D7 Euler/RK4 order --> D8 symplectic/Verlet

 F1 grid --> F2 P1 stiffness --> F3 weak form --> F4 load
                                   |               |
                                   v               v
                               F5 Dirichlet --> F6 order 2
                                   |
                                   v
                              F7 CSR --> F8 CG --> F9 convergence

 M3 d^2=0 --> M6 Stokes --> M7 grad/curl/div  (same root as F3)
\end{verbatim}

\subsection{读法：三层}

\begin{enumerate}[label=（\arabic*）]
  \item \textbf{查漏}：翻到某分支，逐条看「常见误解」——能反驳自己直觉的，才说明真懂；
  \item \textbf{补基础}：条目末的「前置」给出向上追溯的路径，追到 A 开头的条目为止；
  \item \textbf{建直觉}：星号节「串讲」是跨分支的边，专门讲「为什么这两个看起来无关的东西是同一件事」。
\end{enumerate}

本份清单共 \textbf{97 条}知识点：已实现 \textbf{49 条}、框架 \textbf{28 条}、骨架 \textbf{3 条}、
未建 \textbf{16 条}、缺位 \textbf{1 条}。这五类数字的差别比总数重要得多，含义见第 \ref{sec:audit} 节。

状态词是\textbf{白名单五值}、可脚本核对（\texttt{math/survey/verify.py} 的 V2）：
其中「未建」指只在 \texttt{NOW.md} 里规划、「缺位」指库里既无位置也无规划。
两个词不许混用，混了计数就对不上。

\section{分析基础：误差是怎么被量出来的}

这一节是全篇的地基。后面每个阈值、每条「误差 $O(h^p)$」的判据，根都在这里。

\begin{know}{A1}{机器 epsilon 与相对精度}{已实现：core/scalar.hpp}
\field{是什么} 设浮点舍入到最近，$\eps$ 定义为满足 $1+\delta\ne1$ 的最小正 $\delta$。
它给的不是「最小正数」，而是\textbf{相对精度}的量级单位：$\abs{\mathrm{fl}(x)-x}\lesssim \eps\abs{x}/2$。

\field{为什么需要} 全库所有容差都以它为分子；判断「这个残渣算不算零」必须先知道机器能分辨到哪。

\field{关键结论} $\eps_{\text{float}}\approx1.19\times10^{-7}$，$\eps_{\text{double}}\approx2.22\times10^{-16}$，
\textbf{两者差 9 个数量级}。成立条件：舍入到最近；进入次正规数区间后相对误差界失效。

\field{判据} 源码无独立用例，但它是 \texttt{DefaultTol} 的唯一来源。

\field{前置} 无。

\field{常见误解} 以为 $\eps$ 是「最小正浮点数」。那是次正规下界（double 约 $5\times10^{-324}$），
与 $\eps$ 差三百多个数量级；两者混用会让阈值失去量纲意义。
\end{know}

\begin{know}{A2}{累加误差与 $8n\eps$ 的来历}{已实现：core/scalar.hpp}
\field{是什么} $n$ 次量级相当的加减乘，舍入误差累积上界约为 $n\eps\cdot(\text{结果量级})$；
库取系数 $8$ 作余量，定义 $\mathrm{DefaultTol}(n)=8n\eps(\text{real\_t})$。

\field{为什么需要} 它是全库\textbf{唯一默认容差}：判主元是否为零、判 Cholesky 主元是否为正、
判特征值迭代是否收敛，都用它。

\field{关键结论} $8n\eps$ 的量级是「$n$ 次累加的最坏舍入」。成立条件：各步量级相当、
\textbf{没有相消}（见 A6）。一旦出现严重相消，相对误差无界，这个界不再有效。

\field{判据} \texttt{solve.hpp} 判 $\abs{\text{pivot}}\le \mathrm{DefaultTol}(n)$；
\texttt{cholesky.hpp} 判 $\mathrm{real}(d)>\mathrm{DefaultTol}(n)$（写成 \texttt{!(d>tol)} 让 NaN 也抛）。

\field{前置} A1。

\field{常见误解} 以为 $8n\eps$ 是误差分析严格推出来的。源码明确写着「系数 8 是余量、非推导」——
它是个经过实测校准的工程常数，不是定理。
\end{know}

\begin{know}{A3}{前向误差与后向误差}{已实现（概念），但链路未展开}
\field{是什么} 前向误差 $\abs{\hat x-x}/\abs{x}$；后向误差是最小的 $\norm{E}/\norm{A}$，
使 $(A+E)\hat x=b$ 精确成立。前者问「算得准不准」，后者问「算的是不是某个邻近问题的精确解」。

\field{为什么需要} 判断一个算法「稳不稳」必须说清是哪一种。LU 带部分主元是\textbf{后向稳定}的
（$\norm{E}\sim n\eps\norm{A}$），但它的前向误差还要被 $\kappa(A)$ 放大。

\field{关键结论} 一阶近似下 $\text{前向}\lesssim\kappa(A)\cdot\text{后向}$。成立条件：$\kappa$ 取与范数匹配的
2-范数条件数，且扰动不改变秩。

\field{判据} 原文未给——\texttt{solve.hpp} 的「三」段没有展开这条链，属于第 \ref{sec:audit} 节的缺口。

\field{前置} A1、A4。

\field{常见误解} 以为「残差小就等于解准」。残差正是后向误差；$\kappa$ 很大时，
残差 $10^{-16}$ 的解可以离真解差很远。
\end{know}

\begin{know}{A4}{条件数与误差放大}{部分已实现：linalg/svd.hpp 提及}
\field{是什么} $\kappa(A)=\norm{A}\norm{A^{-1}}$；在谱范数下 $\kappa=\sigma_{\max}/\sigma_{\min}$
（对可逆矩阵）。它衡量「输入的小相对扰动最多被放大几倍」。

\field{为什么需要} 解释「为什么解方程的结果比矩阵本身还不准」，也是 SVD 走 $A^{H}A$ 路线
精度损失的唯一解释（N12）。

\field{关键结论} $\kappa(A^{H}A)=\kappa(A)^{2}$；$\sigma$ 的相对误差约 $\sqrt{\eps}$ 而非 $\eps$。
成立条件：$A$ 非奇异；方阵。

\field{判据} 原文只给了一个数字：$\kappa(A)=10^{8}$ 时 $\sigma$ 的相对误差约 $10^{-8}$。
其余文件都「原文未给」（缺口见第 \ref{sec:audit} 节）。

\field{前置} L7。

\field{常见误解} 以为 $\kappa$ 是矩阵的固有性质、与范数无关。换范数 $\kappa$ 会差常数倍；
说 $\kappa$ 必须同时说范数。
\end{know}

\begin{know}{A5}{收敛阶与「h 减半」判据}{已实现（判据方法论）}
\field{是什么} 全局误差 $e(h)=O(h^{p})$ 称 $p$ 阶。实操判据：$h$ 每减半，误差应降 $2^{p}$ 倍，
即 $\log e$ 对 $\log h$ 的斜率约为 $p$。

\field{为什么需要} RK4 的四阶、P1 有限元在 $L^{2}$ 下的二阶，都只能这样验收——没有精确解的场合
这是唯一可用的客观判据。

\field{关键结论} 局部截断误差 $p+1$ 阶 $\Rightarrow$ 全局 $p$ 阶。成立条件：方法稳定、右端项足够光滑、
\textbf{步长已进入渐近区}（$h$ 太大时高阶项未起来，比值会乱）。

\field{判据} 三次 $h$ 减半的比值落在区间内（P1 取 $[3.8,4.2]$）；Euler 减半降 2 倍、RK4 降 16 倍。

\field{前置} A1。

\field{常见误解} 以为「阶越高越准」。阶是渐近率；常数项很大时，低阶方法在中等 $h$ 下可能更准
（RK4 步长过大时反而不如小步长 Euler）。
\end{know}

\begin{know}{A6}{相消与三种改写}{已实现：core/complex.hpp、math/special.hpp}
\field{是什么} 两个相近数相减，结果的相对误差约为 $\eps\abs{a}/\abs{a-b}$——
当 $\abs{a-b}\ll\abs{a}$ 时有效位成批消失。

\field{为什么需要} 四个地方直接踩它：复数模长 $\sqrt{a^{2}+b^{2}}$、开方里的 $\abs{z}\pm a$、
\texttt{erfc} 的大参数、二次方程小根。

\field{关键结论} 三条改写手段：\textbf{换形式}（分子有理化、半角公式）、
\textbf{用不抵消表达式}（\texttt{hypot}、倒数形式）、\textbf{切到互补函数}（$\mathrm{erf}\leftrightarrow\mathrm{erfc}$、
$\Phi(-x)=\mathrm{erfc}(x/\sqrt2)/2$）。成立条件：改写必须在同一数学恒等式的两侧，不能引入新的近似。

\field{判据} \texttt{special.hpp} 要求 $\mathrm{erf}/\mathrm{erfc}$ 与语言内建差 $\le 2$ ulp；
$\mathrm{erfc}$ 在 $x\gtrsim5.9$ 处若用 $1-\mathrm{erf}$ 只剩 1 ulp 可用，必错。

\field{前置} A1。

\field{常见误解} 以为「上 long double 就能治相消」。相消丢的是\textbf{已经算出来的位}，
加精度只多几位，改写形式才根治。
\end{know}

\section{线性代数：矩阵是线性映射，不是数字表}

\subsection*{串讲一：$8n\eps$、回代、条件数是一条链}

A2 给的是「$n$ 步运算能攒多少舍入」；A3 把它翻译成「解方程的后向误差」；
A4 再加上「问题本身的放大倍率」。三者串起来才是完整的误差预算：
\[
\text{解的相对误差}\ \lesssim\ \underbrace{\kappa(A)}_{\text{A4}}\cdot
\underbrace{n\eps}_{\text{A2}}\ .
\]
所以看到 \texttt{DefaultTol(n)}~$=8n\eps$ 时不要只想「机器精度乘 n」，它是这条链的第一环。

\begin{know}{L1}{矩阵作为线性映射与行主序}{已实现：core/matrix.hpp}
\field{是什么} $A\in\R^{m\times n}$ 表示映射 $x\mapsto Ax$（$x$ 为 $n\times1$ 列向量，D19）；
存储上 $A(i,j)$ 落在 \texttt{data}[$i\cdot n+j$]。

\field{为什么需要} 它是所有形状规则与布局规则的源头：向量是列、雅可比是 $m\times n$、
链式法则是 $J_{f\circ g}=J_fJ_g$。

\field{关键结论} 行主序意味着「同一行的元素在内存里相邻」；这是\textbf{存储}约定，
与「向量写成 $n\times1$」不冲突——两者是不同层面的约定。

\field{判据} 无（存储公式本身即是定义）。

\field{前置} 无。

\field{常见误解} 以为行主序 $\Rightarrow$ 向量该写成行向量。布局决定遍历顺序，不决定向量的形状语义。
\end{know}

\begin{know}{L2}{矩阵乘法与形状约束}{已实现：core/matrix.hpp}
\field{是什么} $(AB)_{ij}=\sum_{k=0}^{n-1}A_{ik}B_{kj}$，形状 $(m\times n)(n\times p)=(m\times p)$。

\field{为什么需要} 复合映射就是矩阵乘；雅可比链式法则、反向传播的 $\partial L/\partial W=vx^{\mathsf T}$
都建立在它之上。

\field{关键结论} 中间维必须相等，否则抛异常——库\textbf{不做隐式转置、不做广播}。
成立条件：维数匹配；Hadamard 积走独立入口，不混用。

\field{判据} 无。

\field{前置} L1。

\field{常见误解} 以为逐元素乘（Hadamard）是矩阵乘的特例。两者是不同的运算，
库把它们拆成两个入口正是为了防止混淆。
\end{know}

\begin{know}{L3}{转置、共轭转置与伴随}{已实现：core/matrix.hpp}
\field{是什么} $(A^{\mathsf T})_{ij}=A_{ji}$；$(A^{H})_{ij}=\mathrm{conj}(A_{ji})$。
数学上伴随的定义是 $\inner{A^{H}u}{v}=\inner{u}{Av}$。

\field{为什么需要} 复域一切酉性、正交性、Hermitian 判据都只能用 $A^{H}$；反向传播的转置也是它。

\field{关键结论} 实域 $\mathrm{conj}$ 恒等，故 $A^{\mathsf T}=A^{H}$，\textbf{bug 只在复数域暴露}——
源码记录这个坑已踩两次（\texttt{givens}、\texttt{qr\_inverse}）。换内积（度量矩阵 $G$）时
伴随变成 $G^{-1}A^{H}G$。成立条件：内积为标准 Hermite 内积。

\field{判据} $\norm{Q^{H}Q-I}$、Hermitian 判定，均在 double 下恒真、Complex 下才有效。

\field{前置} C1。

\field{常见误解} 以为「转置 = 改行列布局」。那是实现层面的交换下标；数学上转置是\textbf{伴随}，
与内积绑定，所以换度量它就换形状。
\end{know}

\begin{know}{L4}{内积、共轭位置与正定性}{已实现：linalg/operator.hpp}
\field{是什么} $\inner{a}{b}_{F}=\sum_i \overline{a_i}b_i$（共轭在\textbf{第一个}参数）。

\field{为什么需要} 范数、正交、投影、最小二乘、CG 的正定性前提，全由内积定义。

\field{关键结论} 共轭必须在第一参数，才有 $\inner{x}{x}=\norm{x}^{2}\in\R_{\ge0}$——即\textbf{正定性}。
漏写共轭不会「少一点精度」，而是让内积不再是正定形式，依赖它的正交性判据静默给出错误结论。
成立条件：元素类型具备 $\mathrm{conj}$ 与 $\mathrm{abs2}$ 语义（D7）。

\field{判据} 复数域才测得出（实域共轭恒等，测试免疫）。

\field{前置} C1。

\field{常见误解} 以为「先乘 a 还是先乘 b 无所谓」。复内积是半双线性的，交换要付共轭；顺序错就塌成双线性。
\end{know}

\begin{know}{L5}{Frobenius 范数、迹恒等式与「平方阈值」}{已实现：linalg/operator.hpp}
\field{是什么} $\norm{A}_{F}^{2}=\sum_{ij}\abs{A_{ij}}^{2}$；$\inner{A}{B}_{F}=\mathrm{tr}(A^{H}B)$；
比较一律写成 $\mathrm{abs2}\le \mathrm{tol}^{2}$。

\field{为什么需要} 全库「判大小」的统一写法：避免开方、避免比较复数本身。

\field{关键结论} 平方在非负实数上保序，故 $\abs{x}\le t\iff\abs{x}^{2}\le t^{2}$；
但 $\mathrm{abs2}$ \textbf{丢相位}，只能判模长，判复向量相等必须逐分量比实虚两槽（D7 的已知盲区）。

\field{判据} 所有范数判据（$\norm{A^{\mathsf T}A-I}$、$\norm{A-QR}$、$\mathrm{off}(A)$）都写成平方比较。

\field{前置} L4、A6。

\field{常见误解} 以为 $\mathrm{abs2}$ 判等与 $\abs{\cdot}$ 判等等价。对模长等价；对「两个复数相同」不等价。
\end{know}

\begin{know}{L6}{谱分解与实对称可对角化}{已实现：linalg/eigen.hpp（仅对称/Hermitian）}
\field{是什么} $A=Q\Lambda Q^{H}$，$Q^{H}Q=I$，$\Lambda=\mathrm{diag}(\lambda_1,\dots,\lambda_n)$ 实。

\field{为什么需要} 「旋转坐标系把矩阵变成对角」这件事，是所有对称问题的标准手法；
Jacobi 迭代就是逐步逼近这个 $Q$。

\field{关键结论} 谱定理：$A^{H}=A\Rightarrow$ 正交可对角化、特征值全实、特征向量可取成正交基。
成立条件：\textbf{Hermitian 是结构条件}，不是算法保证——非对称矩阵走这条路径会直接抛错。

\field{判据} 迹不变 $a'+b'=a+b$；$\mathrm{off}(A)\le 8n\eps\norm{A}_{F}$；单元素跳过阈值 $\mathrm{tol}^{2}/n_d$
（保证全部跳过时总和仍不超）。

\field{前置} L3。

\field{常见误解} 以为「实矩阵的特征值都是实的」。只有对称/Hermitian 才保证；一般实矩阵的特征值可成共轭对。
\end{know}

\begin{know}{L7}{奇异值分解与三个恒等式}{已实现：linalg/svd.hpp}
\field{是什么} $A=USV^{H}$，$U^{H}U=I_r$、$V^{H}V=I_r$，$r=\min(m,n)$，
$\sigma_i=\sqrt{\lambda_i(A^{H}A)}\ge0$ 且降序。

\field{为什么需要} 秩、条件数、最小二乘、低秩逼近都由它给出；它也是唯一能在数值上「看清」$\kappa$ 的工具。

\field{关键结论} 三条恒等式：$\sum_i\sigma_i^{2}=\norm{A}_{F}^{2}$（Frobenius）、
$\sum_i\sigma_i=\norm{A}_{*}$（核范数）、$\sigma_{\max}\sigma_{\min}=\abs{\det A}$（方阵非奇异）。
左奇异向量 $u_i=Av_i/\sigma_i$，$\sigma_i\approx0$ 时 $u_i$ 不唯一（可任取正交补中的向量）。

\field{判据} 秩用相对阈值：$\mathrm{rank}(A)=\#\{i:\sigma_i>\tau\sigma_0\}$，$\tau=\max(m,n)\eps$；
奇异值判零用的 $\mathrm{tol\_sig}=\mathrm{DefaultTol}(r)\cdot\sigma_0$。
（任务说明里提到的「可信阈值 $1\times10^{-6}$」\textbf{原文未给}。）

\field{前置} L6。

\field{常见误解} 以为 $\sigma_i>0$ 就能判秩。浮点下 $\lambda$ 可能出现 $-10^{-17}$，
必须 $\max(\lambda,0)$ 夹取并用相对阈值；绝对阈值在尺度变化时全错。
\end{know}

\begin{know}{L8}{正交相似与不变量}{已实现：linalg/eigen.hpp}
\field{是什么} $A'=G^{H}AG$，$G$ 酉（$G^{H}G=I$）。

\field{为什么需要} Jacobi 每做一次旋转都在做这个变换，靠不变量判断「有没有把矩阵搞坏」。

\field{关键结论} 酉相似保持：谱、Frobenius 范数、迹、Hermitian 性。成立条件：$G$ 必须\textbf{酉}；
非正交相似只保谱，会把范数放大。

\field{判据} 「迹不变」是唯一廉价可查的判据（$a'+b'=a+b$）；$\norm{A}_{F}$ 只算一次，后面对拍。

\field{前置} L3、L5。

\field{常见误解} 以为「相似就保范数」。$G$ 不正交时 $\norm{G^{H}AG}$ 可以任意放大，迭代的收敛判据随之失效。
\end{know}

\begin{know}{L9}{行列式是体积，不是「接近 0 就奇异」}{库无 det 入口（原文未给）}
\field{是什么} $\det A$ 是有向体积缩放因子；方阵下 $\abs{\det A}=\prod_i\sigma_i$。

\field{为什么需要} 判断坐标变换的 Jacobian 是否为零（拉回与积分换元要它），以及体积判据。

\field{关键结论} $\abs{\det A}$ 小\textbf{不等于}病态：把单位矩阵乘 $10^{-8}$，$\det=10^{-8n}$ 而 $\kappa=1$。
判断奇异性要用 $\sigma_{\min}/\sigma_{\max}$。

\field{判据} 原文未给——库没有 \texttt{det} 入口，这条属缺口。

\field{前置} L7。

\field{常见误解} 以为 $\det\approx0$ 就是奇异。尺度一改判断就翻，它是个\textbf{绝对量}，不能横向比较。
\end{know}

\begin{know}{L10}{秩 1 修正（外积）}{已实现：linalg/qr.hpp、autodiff/grad\_graph.hpp}
\field{是什么} $uv^{H}$ 是秩 1 矩阵，$(uv^{H})_{ij}=u_i\overline{v_j}$。

\field{为什么需要} Householder 反射 $H=I-2vv^{H}/(v^{H}v)$ 就是秩 1 修正；
反向传播里 $\partial L/\partial W=vx^{\mathsf T}$ 也是外积。

\field{关键结论} 秩 1 修正只改变一个方向上的分量，故一次反射能把一个向量整体打到坐标轴上。
成立条件：$u,v$ 为列向量（D19）；形状必须匹配，否则抛错。

\field{判据} 反向模式的形状校验：$\partial L/\partial W$ 必须是 $m\times n$，由 $vx^{\mathsf T}$ 唯一确定（G4）。

\field{前置} L2、L3。

\field{常见误解} 以为外积是内积的「转置版」。内积给标量、外积给矩阵，维数差一个量级。
\end{know}

\section{数值线性代数：分解是为了让「判据」有地方落}

\begin{know}{N1}{部分主元的必要性}{已实现：linalg/solve.hpp}
\field{是什么} 消元第 $k$ 步前，在当前列的第 $k$ 行以下选绝对值最大的元素换到对角位置，
再做消元。换行信息记在 \texttt{piv} 里。

\field{为什么需要} 不选主元时，主元可能是 $10^{-20}$，下一步 $1/10^{-20}=10^{20}$ 会把残差放大到
完全不可用的量级——这\textbf{不是精度问题，是算法失效}。

\field{关键结论} 选主元保证 $\abs{l_{ik}}\le1$（乘子不放大），从而 $L$ 的元素有界、
后向误差 $\norm{E}\sim n\eps\norm{A}$。成立条件：无结构假设；对 SPD 矩阵可省（见 N7）。

\field{判据} 主元判零用 $\mathrm{abs2}(p)\le\mathrm{tol}^{2}$，$\mathrm{tol}=\mathrm{DefaultTol}(n)=8n\eps$。

\field{前置} A2。

\field{常见误解} 以为选主元是为了「更准」。它首先是为了\textbf{不炸}；精度提升是副产品。
\end{know}

\begin{know}{N2}{LU 分解与乘子界}{已实现：linalg/solve.hpp}
\field{是什么} $PA=LU$：$L$ 单位下三角（对角全 1），$U$ 上三角，$P$ 由行交换序列重放。
消元式：$l_{ik}=u_{ik}/u_{kk}$，$u_{ij}\leftarrow u_{ij}-l_{ik}u_{kj}$（$j\ge k$）。

\field{为什么需要} 一次分解、多次求解：同一 $A$ 换不同 $b$ 只需前代回代，代价从 $O(n^{3})$ 降到 $O(n^{2})$。

\field{关键结论} 前代 $Ly=Pb$、回代 $Ux=y$；不显式求 $A^{-1}$（求逆的 flops 是解方程的数倍，精度还更差）。
成立条件：$A$ 方阵、主元非奇异；$\abs{l_{ij}}\le1$ 由 N1 保证。

\field{判据} 主元 $\abs{u_{kk}}>\mathrm{DefaultTol}(n)$，否则抛「输入合法但数值做不到」类异常。

\field{前置} N1。

\field{常见误解} 以为分解完直接调逆矩阵更方便。$A^{-1}$ 的每个元素都要解一次方程，
而且残差更大——凡是只要 $x$，就别造 $A^{-1}$。
\end{know}

\begin{know}{N3}{Householder 反射}{已实现：linalg/qr.hpp}
\field{是什么} 取 $v=x+\mathrm{sign}(x_1)\norm{x}e_1$，令 $H=I-2vv^{H}/(v^{H}v)$，则
$Hx=-\mathrm{sign}(x_1)\norm{x}e_1$——一次反射把一个向量整体打到坐标轴上。

\field{为什么需要} 它是 QR 的「整列消元」工具：一次反射消掉一列下面的所有元素，不改已算好的列。

\field{关键结论} $H^{H}H=I$ 且 $H^{H}=H$（酉且 Hermitian，故 $H^{-1}=H$）。
分母 $v^{H}v$ 必须用\textbf{修改后}的 $v$ 计算——用改前的 $\norm{x}^{2}$ 会让 $H$ 不再对合，
消元目标也达不到。成立条件：$x\ne0$；$\norm{x}$ 大于阈值，否则跳过。

\field{判据} 跳过条件 $\norm{x}\le \mathrm{DefaultTol}(m)$。

\field{前置} L3、L4、L10。

\field{常见误解} 以为反射会放大误差。酉变换的 2-范数恰为 1，它是「不放大」的；
真正放大误差的是后续的消元除法。
\end{know}

\begin{know}{N4}{Givens 旋转}{已实现：linalg/qr.hpp}
\field{是什么} $G=\begin{bmatrix}c&s\\-\overline{s}&c\end{bmatrix}$，$\abs{c}^{2}+\abs{s}^{2}=1$；
取 $r=\mathrm{hypot}(\abs{f},\abs{g})$，$c=\abs{f}/r$，$s=(f/\abs{f})\cdot\overline{g}/r$，
则 $G(f,g)^{\mathsf T}=((f/\abs{f})r,0)^{\mathsf T}$。

\field{为什么需要} 一次只消一个元素，适合\textbf{稀疏}或「只需局部消元」的场合；
反射适合稠密整列。

\field{关键结论} 复数域里靠相位 $f/\abs{f}$ 补一个自由度，才能让 $G$ 保持酉且把 $g$ 精确消成 0。
成立条件：$\abs{f}$、$\abs{g}$ 不同时为 0；用 \texttt{hypot} 防上溢。

\field{判据} $\abs{g}^{2}\le\mathrm{tol}^{2}$ 时跳过（已足够小）。

\field{前置} L3、L4。

\field{常见误解} 以为 Givens 只是「2×2 的反射」。它是旋转（行列式为 1），不是反射（行列式 $-1$）；
差别在复数域与相位处理上会现形。
\end{know}

\begin{know}{N5}{MGS 与 CGS：病态下才见分晓}{已实现：linalg/qr.hpp（MGS）}
\field{是什么} 经典 Gram--Schmidt（CGS）先算出所有 $r_{ij}=q_i^{H}a_j$ 再一次性扣；
修正版（MGS）每算一个系数就立刻从 $a_j$ 扣掉。

\field{为什么需要} 它是最直观的正交化；但在病态矩阵上 CGS 的正交性会垮掉，必须知道为什么。

\field{关键结论} MGS 的舍入误差约 $\eps\kappa(A)$，CGS 约 $\eps\kappa(A)^{2}$。
源码记录：CGS 在病态下 $\norm{Q^{H}Q-I}$ 可到 $1\times10^{-8}$（double 下这已是灾难）。
成立条件：$q_k$ 必须已归一（$\inner{q_k}{q_k}\equiv1$，分母被吸收）。

\field{判据} $\norm{A-QR}_{\infty}$ 与 $\norm{Q^{H}Q-I_{n}}_{\infty}$ 同时小；
秩亏列 $r_{jj}=0$ 时 $Q$ 的第 $j$ 列不唯一，不能用「逐列一致」当判据。

\field{前置} L4、A4。

\field{常见误解} 以为 MGS 与 CGS 只是「循环顺序不同、结果一样」。数学上等价，浮点上\textbf{精度差一个 $\kappa$ 量级}。
\end{know}

\begin{know}{N6}{thin QR 与最小二乘}{已实现（QR），最小二乘原文未给}
\field{是什么} $A=QR$：$A$ 为 $m\times n$（$m\ge n$），$Q$ 为 $m\times n$ 且 $Q^{H}Q=I_n$，
$R$ 为 $n\times n$ 上三角（thin QR）。

\field{为什么需要} 它是\textbf{最小二乘}的标准路径：$\min\norm{Ax-b}_{2}$ 等价于解 $Rx=Q^{H}b$，
因为正交变换不改变 2-范数。

\field{关键结论} 最小二乘的数学链是
\[
\norm{Ax-b}^{2}=\norm{Q^{H}(Ax-b)}^{2}
=\norm{Rx-Q^{H}b}^{2}+\norm{(I-QQ^{H})b}^{2},
\]
第二项与 $x$ 无关，故只需解上三角方程。成立条件：$A$ 列满秩（否则 $R$ 奇异，需换 QR 带列选主元或 SVD）。

\field{判据} 原文未给——\texttt{qr.hpp} 的「三」段\textbf{没有写最小二乘}，这条属缺口（见第 \ref{sec:audit} 节）。
$m<n$ 的情形库统一抛错，不做退化处理。

\field{前置} N3。

\field{常见误解} 以为 $Q$ 必须是 $m\times m$ 的方阵。thin QR 只取前 $n$ 列，
这正是「经济型」的含义；补齐多余列只改变 $Q$，不改 $R$。
\end{know}

\begin{know}{N7}{Cholesky：对称正定的专用分解}{已实现：linalg/cholesky.hpp}
\field{是什么} $A=LL^{H}$，$L$ 下三角且 $L(j,j)>0$（取正号即唯一）。
逐列：$d=A(j,j)-\sum_{k<j}\abs{L(j,k)}^{2}$，$L(j,j)=\sqrt d$，
$L(i,j)=\big(A(i,j)-\sum_{k<j}L(i,k)\overline{L(j,k)}\big)/L(j,j)$。

\field{为什么需要} SPD 系统的标准解法：代价约 $n^{3}/3$（LU 是 $2n^{3}/3$），且\textbf{不需要选主元}——
正定性保证主元严格为正。

\field{关键结论} 唯一性来自「对角取正」这一条约定；若允许负号，$L$ 有 $2^{n}$ 种。
成立条件：$A$ Hermitian 正定；只读下半三角，故非 Hermitian 输入不会报错（由调用方保证）。

\field{判据} $\mathrm{real}(d)>\mathrm{DefaultTol}(n)$，写成 \texttt{!(d>tol)} 让 NaN 也落入失败分支。

\field{前置} L4、A2。

\field{常见误解} 以为「正定」可以用「所有对角元为正」来判。反例：$\begin{bmatrix}1&2\\2&1\end{bmatrix}$
对角全正但不定；判据只有分解走通（N8）。
\end{know}

\begin{know}{N8}{SPD 判定 = 分解走通}{已实现：inverse.hpp、cholesky.hpp}
\field{是什么} $A$ Hermitian 正定 $\iff$ 不选主元的 Cholesky 能完整走通。

\field{为什么需要} 这条等价性让「判定」和「计算」合并成一步——不需要单独的判定算法，
跑一遍分解就是判定。

\field{关键结论} 双向等价：正定 $\Rightarrow$ 每步主元 $d>0$；反之若每步都走通，则构造出的 $L$ 满足 $LL^{H}=A$
且主元为正 $\Rightarrow$ 正定。成立条件：$A$ 必须 Hermitian（否则「正定」本身无定义）。

\field{判据} 主元 $\mathrm{real}(d)>\mathrm{DefaultTol}(n)$，失败即判非 SPD（够用，因为主元严格正时有正的下界）。

\field{前置} N7。

\field{常见误解} 以为 SPD 判定需要先算特征值。特征值判定正确但要 $O(n^{3})$ 且精度更差；
分解判定是同一量级里最省的路线。
\end{know}

\begin{know}{N9}{结构化解与探测顺序}{已实现：linalg/inverse.hpp}
\field{是什么} 按结构给出逆：对角 $(A^{-1})_{ii}=1/d_i$；正交 $A^{-1}=A^{H}$；
SPD 用 $A=LL^{H}\Rightarrow A^{-1}=(L^{-1})^{H}L^{-1}$；一般用 $PA=LU\Rightarrow A^{-1}=U^{-1}L^{-1}P$；
QR 用 $A^{-1}=R^{-1}Q^{H}$。

\field{为什么需要} 结构化解又快又准（正交矩阵的逆零误差），比通用求解好得多；但前提是\textbf{判对结构}。

\field{关键结论} 结构探测的集合关系：对角 $\subset$ 对称 $\subset$ 一般；SPD $\subset$ 对称。
所以在检查顺序上\textbf{子类必须判在父类之前}，否则永远走通用路径。
条件判据：对角 $\abs{a_{ij}}\le\mathrm{tol}\,(i\ne j)$；对称 $\abs{a_{ij}-\overline{a_{ji}}}\le\mathrm{tol}$；
正交 $\norm{A^{H}A-I}\le\mathrm{tol}$；复数域必须用 $A^{H}$ 而不是 $A^{\mathsf T}$。

\field{判据} 各路径的奇异判据统一：对角 $\abs{d_i}\le\mathrm{tol}$，LU $\abs{u_{kk}}\le\mathrm{tol}$，
QR $\abs{R_{ii}}\le\mathrm{tol}$，SPD 由 N8 内部判。

\field{前置} N2、N7、N8。

\field{常见误解} 以为「对称矩阵的逆也对称」就能反推用了快路径。逆确实对称，
但判定它是否对称用的容差是调用方给的——同一个矩阵在 $10^{-6}$ 与 $10^{-12}$ 下可能落到不同分支。
\end{know}

\begin{know}{N10}{矩阵指数：Padé 与 scaling \& squaring}{已实现：linalg/expm.hpp}
\field{是什么} 两条路线：Taylor 级数截断；Padé $[13/13]$ 近似
$\exp(A)\approx D^{-1}N$，其中 $N=\sum_{k=0}^{m}c_kA^{k}$、$D=\sum_k c_k(-A)^{k}$，
$c_k=\frac{(2m-k)!\,m!}{(2m)!\,k!\,(m-k)!}$（递推 $c_k=c_{k-1}\frac{m-k+1}{k(2m-k+1)}$，$c_1=1/2$）。
两者都配 scaling \& squaring：$\exp(A)=(\exp(A/2^{s}))^{2^{s}}$。

\field{为什么需要} 矩阵指数是线性 ODE、旋转积分、马尔可夫半群的统一工具；
它也是「逼近算法 + 误差地板」这一对概念的样板。

\field{关键结论} 取 $s$ 使 $\norm{A/2^{s}}_{F}\le1/2$（Taylor）或 $\le\theta_{13}=5.3719\ldots$（Padé）。
\textbf{精度地板}：$\exp$ 的条件数约 $e^{\norm{A}}/\sqrt{2\pi\norm{A}}\cdot\eps$，与算法无关——
即「再怎么改算法，相对误差也下不去」。
独立判据：$\det(\exp A)=\exp(\mathrm{tr}A)$；而 $\exp(A)\exp(-A)=I$ 在 $\norm{A}=100$ 时条件数可达 $7\times10^{86}$，
不能当判据。

\field{判据} Padé 与 Taylor 的相对差 $<10^{-8}$（实测 $5.3\times10^{-16}$ 至 $1.5\times10^{-13}$）；
级数截断上界 $k\le30$。

\field{前置} N2、A4。

\field{常见误解} 以为「$\exp(A)$ 就是把每个元素取 $\exp$」。那是逐元素操作，不是矩阵指数；
只有对角矩阵两者才重合。
\end{know}

\begin{know}{N11}{Jacobi 旋转求对称特征值}{已实现：linalg/eigen.hpp}
\field{是什么} 对 $2\times2$ Hermitian 子块 $\begin{bmatrix}a&z\\\overline{z}&b\end{bmatrix}$，
用酉旋转零化 $z$：$\tan2\theta=\frac{2\abs{z}}{a-b}$，$\theta=\frac12\mathrm{atan2}(2\abs{z},a-b)$；
旋转后 $a'=\frac{a+b}{2}+\frac{D}{2}$，$b'=\frac{a+b}{2}-\frac{D}{2}$，$D=\mathrm{hypot}(a-b,2\abs{z})$。

\field{为什么需要} 它是最稳的对称特征值算法：只做正交相似（L8），不破坏结构，收敛二次。

\field{关键结论} 对角元用解析式而不写 $c^{2}a+s^{2}b+2cs\abs{z}$——后者多一次相消（A6）。
$\theta$ 若取成 $\theta+\pi/2$ 会把两特征值对调，故公式与取法必须配套。
收敛判据必须\textbf{相对}：$\mathrm{off}(A)\le\mathrm{tol}\norm{A}_{F}$，
其中 $\mathrm{off}(A)^{2}=\norm{A}_{F}^{2}-\sum_i\abs{A_{ii}}^{2}$；绝对阈值在 $\norm{A}$ 大时永不收敛、
小时过早停。成立条件：$A$ Hermitian；每轮 cyclic 扫全部 $i<j$（否则会填回已零化的元素）。

\field{判据} 单元素跳过阈值 $\mathrm{tol}^{2}/n_d$，$n_d=n(n-1)$；最多 $kMaxSweep=100$ 轮，
超限抛异常；迹不变可随时对拍。

\field{前置} L6、L8、N3。

\field{常见误解} 以为「零化一个元素就少一个非对角元」。旋转会把别处已零化的元素重新填成非零，
所以必须扫完整轮并看整体 $\mathrm{off}$。
\end{know}

\begin{know}{N12}{为什么走 $A^{H}A$ 会把条件数平方}{已实现：linalg/svd.hpp（路线选择）}
\field{是什么} 由于 $x^{H}A^{H}Ax=\norm{Ax}^{2}\ge0$，$B=A^{H}A$ 是 Hermitian 半正定，
其特征值就是 $\sigma_i^{2}$；于是 SVD 可以退化为「对称特征值问题」。

\field{为什么需要} 这是库当前 SVD 的实现路线——最省事、复用已有的 Jacobi；但代价必须说清。

\field{关键结论} $\kappa(A^{H}A)=\kappa(A)^{2}$（这是 A4 的直接推论）。因此在 $B$ 上求特征值的相对误差
$\sim\eps\kappa(B)=\eps\kappa(A)^{2}$，开方后 $\sigma$ 的相对误差 $\sim\sqrt{\eps}\,\kappa(A)$。
源码给的例子：$\kappa(A)=10^{8}$ 时 $\sigma$ 相对误差约 $1\times10^{-8}$。成立条件：$\kappa$ 用 2-范数；
重构 $\norm{USV^{H}-A}$ 仍是 $\eps$ 级（$\sigma$ 与 $1/\sigma$ 精确相消）。

\field{判据} $\mathrm{tol\_sig}=\mathrm{DefaultTol}(r)\sigma_0$；补列时要求投影后 $\norm{w}^{2}>0.25$。

\field{前置} A4、L7、N11。

\field{常见误解} 以为「既然数学上等价，走 $A^{H}A$ 没损失」。等价性对精确算术成立；
浮点上损失的不是 $\eps$ 而是 $\sqrt{\eps}$——这是路线选择问题，不是实现 bug。
\end{know}

\begin{know}{N13}{秩判定与阈值的尺度}{已实现：linalg/svd.hpp}
\field{是什么} $\mathrm{rank}(A)=\#\{i:\sigma_i>\tau\sigma_0\}$，$\tau=\max(m,n)\eps$。

\field{为什么需要} 秩是「这个矩阵的信息还剩多少」的量化；没有相对阈值，秩随尺度变化，
判定毫无意义。

\field{关键结论} 阈值必须\textbf{相对最大奇异值}，不能用 $\sigma_i>0$，也不能用固定常数。
成立条件：$\sigma_0>0$（非零矩阵）。$8n\eps$ 型阈值适合判「唯一性/失败」，\textbf{不适合判秩}——
源码记录过用 $8n\eps$ 判秩只对 $73/100$。

\field{判据} $\tau=\max(m,n)\eps$；同时用 $\mathrm{DefaultTol}(r)\sigma_0$ 判「该不该继续算左奇异向量」。

\field{前置} L7、A4。

\field{常见误解} 以为「秩亏」是一个\textbf{是/否}的事实。数值上它是连续量：$\sigma_{\min}$ 从 $10^{-16}$ 连续变到 $10^{-2}$，
秩只在给定阈值下才有定义。
\end{know}

\subsection*{串讲二：正交化三兄弟——同一条岭的三条路}

Householder（N3）、Givens（N4）、MGS（N5）都在做「把向量组变正交」，差别是\textbf{每次处理多少}：

\begin{center}
\begin{tabular}{L{2.1cm}L{2.8cm}L{3.9cm}L{4.8cm}}
\toprule
方法 & 一次处理的量 & 代价 & 稳定性来源 \\
\midrule
Householder & 一整列 & $2n^{2}m-\frac23n^{3}$ 级 & 酉变换，2-范数不放大 \\
Givens & 一列中的一个元素 & 约反射的 2--3 倍 & 同上，但单位旋转更细 \\
MGS & 一个向量对前 $j-1$ 个方向 & 约 $2mn^{2}$ & 只靠「边算边扣」抑制误差，无酉性保证 \\
\bottomrule
\end{tabular}
\end{center}

三者的稳定的本质差别：前两个是\textbf{酉变换}（$H^{H}H=I$，$\norm{Hx}=\norm{x}$ 精确成立），
MGS 只是一次次改写，误差按 $\eps\kappa(A)$ 增长而不是 $\eps\kappa^{2}(A)$。
所以「稠密整列用反射、稀疏局部用旋转、图省事用 MGS 但要接受精度打折」是全部结论。

\section{复分析：多值性不是缺陷，是定义的一部分}

\begin{know}{C1}{复代数与非有序域}{已实现：core/complex.hpp}
\field{是什么} $(a+bi)(c+di)=(ac-bd)+(ad+bc)i$，$\abs{z}^{2}=a^{2}+b^{2}$，
$z^{-1}=\overline{z}/\abs{z}^{2}$。

\field{为什么需要} 它是全库唯一的「非有序标量」：所有比较必须绕道 $\mathrm{abs2}$（D7），
所有正交必须用共轭转置（L3）。

\field{关键结论} 复数域\textbf{不存在}与加法乘法相容的全序。成立条件：无——
这是定理（任何把 $\C$ 变成有序域的尝试都会导致 $-1\ge0$）。

\field{判据} 判复数相等必须逐分量比实虚两槽；$\mathrm{abs2}$ 只能判模长（L5）。

\field{前置} L4。

\field{常见误解} 以为「复数也能比大小，只是编译器不给」。不是编译器限制，是数学上不可能：
只要 $\abs{\cdot}$ 与加减乘相容，$i^{2}=-1<0$ 必然矛盾。
\end{know}

\begin{know}{C2}{Smith 除法}{已实现：core/complex.hpp}
\field{是什么} 朴素式 $(a+bi)/(c+di)=\big[(ac+bd)+(bc-ad)i\big]/(c^{2}+d^{2})$ 会在 $c,d$ 极大时上溢、
极小时下溢。Smith 形式：若 $\abs{c}\ge\abs{d}$，取 $r=d/c$，
则商 $=\dfrac{a+br}{c+dr}+\dfrac{b-ar}{c+dr}i$；$\abs{c}<\abs{d}$ 时对称地按 $d$ 归一化。

\field{为什么需要} 中间量必须落在可表示范围：$c^{2}+d^{2}$ 在 $\abs{c},\abs{d}\sim10^{200}$ 时直接溢出成
$\infty$，商变成 $0$ 或 NaN，而真实商可能就在 1 附近（如 $(10^{200}+i)/(10^{200}+i)=1$）。

\field{关键结论} 先按大分量归一化，使 $\abs{r}\le1$，所有中间量不超过约 $2\max(\abs{c},\abs{d})$。
成立条件：分母非零（零由 IEEE 传播，不在标量层判，见 D12）。

\field{判据} 用例：同量级大数相除应为 $1$；$\abs{c},\abs{d}\sim10^{-200}$ 时不出现 inf/nan。

\field{前置} C1、A6。

\field{常见误解} 以为「先算模方再除」只是慢一点。它是\textbf{错}：模方的动态范围是原分量的两倍，
越界即坏，且无法靠提高精度根治。
\end{know}

\begin{know}{C3}{欧拉公式与复指数}{已实现：core/complex.hpp、linalg/expm.hpp 对照}
\field{是什么} $\exp(a+bi)=e^{a}(\cos b+i\sin b)$。

\field{为什么需要} 它是「旋转」的解析表达：单位圆上的 $\exp(i\theta)$ 就是转动 $\theta$；
DFT 的旋转因子、特征函数的 $e^{itX}$、四元数的半角公式，都是它的变体。

\field{关键结论} 它与矩阵指数 $\exp(A)$ \textbf{只在 $1\times1$ 时重合}。成立条件：$b$ 为实数；
相位按模 $2\pi$ 有意义。

\field{判据} 无（恒等式本身）；\texttt{expm.hpp} 明确区分两者（N10）。

\field{前置} C1。

\field{常见误解} 以为 $\exp(A)$ 可以先对元素取 $\exp$。见 N10；两者的「指数」含义不同：
一个是标量函数，一个是幂级数的矩阵和。
\end{know}

\begin{know}{C4}{对数主值与分支切割}{已实现：core/complex.hpp}
\field{是什么} $\log z=\ln\abs{z}+i\arg z$，$\arg z\in(-\pi,\pi]$（主值）。

\field{为什么需要} 复对数是多值（$\log z$ 差 $2\pi ik$），必须选一支；选哪一支决定了哪些地方不连续。

\field{关键结论} 主值分支的\textbf{切割在负实轴}：从上方趋近给 $+\pi$，从下方给 $-\pi$，
所以 $\log$ 在负实轴上不连续。这是多值性的必然，不是实现缺陷。需要连续性时要自己做相位展开（unwrapping）。

\field{判据} 无。

\field{前置} C3。

\field{常见误解} 以为「主值就是唯一正确的值」。主值是\textbf{约定}；换一支（比如 $[0,2\pi)$）
得到另一套连续性与另一套公式，物理结果相同但中间步骤的符号会变。
\end{know}

\begin{know}{C5}{开方主值与半角写法}{已实现：core/complex.hpp}
\field{是什么} $\sqrt{z}=\sqrt{\dfrac{\abs{z}+a}{2}}+\mathrm{sign}(b)\sqrt{\dfrac{\abs{z}-a}{2}}\,i$，
其中 $a=\mathrm{Re}\,z$。

\field{为什么需要} 直接算 $\sqrt{a+bi}$ 需要解二次方程，会遇相消（A6）；半角形式把它变成
两个非负量之和。

\field{关键结论} 实部恒取正号 $\Rightarrow$ 结果的实部 $\ge0$（这是主值的定义）；
$z=0$ 必须单独返回，否则 $b=0/0$ 变 NaN；跨负实轴时虚部符号翻转，故\textbf{不连续}。
$\abs{z}$ 要用 \texttt{hypot}，$\abs{z}\pm a$ 要用不抵消的写法。成立条件：$z\ne0$ 走通用分支。

\field{判据} $\abs{a},\abs{b}\sim10^{200}$ 时不能出现 inf/nan；$z<0$ 时结果应为纯虚且符号与 $b$ 一致。

\field{前置} C4、A6。

\field{常见误解} 以为 $\sqrt{z^{2}}=\abs{z}$。复域里 $\sqrt{z^{2}}$ 依赖分支，可能是 $z$ 也可能是 $-z$。
\end{know}

\begin{know}{C6}{幂主值与「不是实数幂」}{已实现：core/complex.hpp}
\field{是什么} $z^{w}=\exp(w\log z)$，取 $\log$ 的主值即得幂的主值。

\field{为什么需要} 它说明「分数次幂」在复域天然多值；求根、算 $\sqrt[3]{\cdot}$ 都要先决定分支。

\field{关键结论} $(-8)^{1/3}$ 的主值是 $1+1.732i$，\textbf{不是} $-2$。成立条件：$z\ne0$（$0^{0}$ 与 $0^{w}$ 另议）。

\field{判据} 无。

\field{前置} C4、C5。

\field{常见误解} 以为复数的幂可以「先化成实数再算」。一旦底数为负或指数非整数，
实数路径给出的那个根只是若干分支中的一个，且常不是主值。
\end{know}

\section{调和分析：变换的两端与归一化的位置}

\begin{know}{H1}{DFT 的定义与「正向不归一化」}{框架：math/fourier.hpp}
\field{是什么} $X_k=\sum_{n=0}^{N-1}x_n e^{-2\pi ikn/N}$，$k=0,\dots,N-1$；正向\textbf{不除 $N$}。

\field{为什么需要} 它是「把一个离散信号分解成 $N$ 个旋转的叠加」；不归一化让正向成为纯投影，
所有归一化集中到逆变换，避免因子分散在两处。

\field{关键结论} 旋转因子 $e^{-2\pi ikn/N}$ 对应 C3 的复指数；$X_k$ 是 $x$ 在频率 $k$ 上的\textbf{内积}
$\inner{e_k}{x}$（严格说要共轭，方向由符号约定决定）。成立条件：$N\ge1$；信号按 $N\times1$ 列向量（D19）。

\field{判据} 与朴素 DFT 对拍：$N=1024$ 复数随机，$\norm{\mathrm{FFT}-\mathrm{DFT}}_{\infty}/\norm{X}_{\infty}\le10^{-13}$；
另加\textbf{有条件逐点}判据——只对 $\abs{X_k}\ge0.1\norm{X}_{\infty}$ 的频点比较（近零频点必然假红）。

\field{前置} C3、L4。

\field{常见误解} 以为「不归一化」是无所谓的选择。它决定了卷积定理里有没有附加因子（H3）、
Parseval 里 $1/N$ 放在哪边（H8）——一处改，处处变。
\end{know}

\begin{know}{H2}{逆变换：$1/N$ 只出现一次}{框架：math/fourier.hpp}
\field{是什么} $x_n=\frac1N\sum_k X_k e^{+2\pi ikn/N}$。

\field{为什么需要} 正逆必须复合为恒等：$\frac1N\sum_k e^{2\pi ik(n-m)/N}=\delta_{nm}$。
这个 $1/N$ 是「$N$ 个正交方向的平均」的代价。

\field{关键结论} 两个方向的\textbf{符号相反}是定义本身，不是笔误；$1/N$ 全库只在 IFFT 出现一次。
成立条件：$N$ 与正向一致；复数域共轭关系 $x=\overline{\mathrm{FFT}(\overline{X})}/N$ 可作为实现路径。

\field{判据} 往返 $\norm{\mathrm{IFFT}(\mathrm{FFT}(x))-x}_{\infty}\le10^{-13}$。

\field{前置} H1。

\field{常见误解} 以为「正变换除 $\sqrt N$、逆变换除 $\sqrt N$」更对称。那是酉归一化（Parseval 无因子），
与 D20 约定的 $\pm1$、$1/N$ 方案不同；混用会让卷积少（或多）一个因子。
\end{know}

\begin{know}{H3}{循环卷积定理}{框架：math/fourier.hpp}
\field{是什么} $(a\circledast b)_n=\sum_{m}a_m b_{(n-m)\bmod N}$，则
$\mathrm{FFT}(a\circledast b)_k=A_kB_k$（\textbf{无附加因子}）。

\field{为什么需要} 时域卷积变成频域乘法，是 FFT 全部实用价值的来源。

\field{关键结论} 前提是\textbf{循环}卷积（下标模 $N$）；无附加因子正是不归一化 + $1/N$ 只在逆变换的后果。
成立条件：两序列等长 $N$；不等长要先补零到同一长度，否则得到的是循环的、绕回的结果。

\field{判据} 三条路径对拍：时域直接、频域补零、循环矩阵乘，逐点 $\le10^{-12}$。

\field{前置} H1、H2。

\field{常见误解} 以为「FFT 卷一下就得到线性卷积」。循环卷积会把尾部\textbf{绕回}头部；
结果看起来「差不多」但边界全错，这是最常见的静默错。
\end{know}

\begin{know}{H4}{线性卷积与补零长度}{框架：math/fourier.hpp}
\field{是什么} $(a*b)_n=\sum_m a_mb_{n-m}$，长度 $N+M-1$；用 FFT 时补零到
$L\ge N+M-1$ 且 $L=2^{k}$。

\field{为什么需要} 线性卷积是滤波、多项式乘法的标准形式；只有补够零，循环卷积才与它相等。

\field{关键结论} $L<N+M-1$ 时尾部环绕（$y_{n}$ 与 $y_{n+L}$ 混在一起），这是\textbf{另一个函数}而非精度问题。
源码取 $L=\text{NextPowerOfTwo}(N+M-1)$，且\textbf{不做自动分派}——非 2 的幂不会偷偷退回 $O(N^{2})$ 的 DFT。

\field{判据} 前 $N+M-1$ 项与直接卷积逐点 $\le10^{-12}$；$L$ 取小一位时应当出现可见差异（反例用例）。

\field{前置} H3。

\field{常见误解} 以为「补零只是形式」。补零长度决定\textbf{语义}：够长才是线性卷积，
不够长就是循环卷积，两者在边界差一整段。
\end{know}

\begin{know}{H5}{频轴与 Nyquist}{框架：math/fourier.hpp}
\field{是什么} 采样间隔 $dx$ 下，$\xi_k=k/(N\,dx)$（$k\le N/2$）；后半段
$\xi_k=(k-N)/(N\,dx)$（$k>N/2$）。采样率 $f_s=1/dx$，Nyquist 频率 $f_s/2$ 落在 $k=N/2$。

\field{为什么需要} 没有频轴，$X_k$ 只是「第 $k$ 个数」；有了频轴才能说「这是 3 Hz 的分量」。

\field{关键结论} 前半段是正频率、后半段是负频率（因为 DFT 是周期谱，$\xi$ 与 $\xi+f_s$ 不可分辨）。
成立条件：等距采样；$dx>0$。$N$ 偶数时 $k=N/2$ 是自共轭点（实信号下该点相位移 $\pi$ 无意义）。

\field{判据} 直流分量在 $k=0$；实信号在 $k$ 与 $N-k$ 共轭（H8 的对称性）。

\field{前置} H1。

\field{常见误解} 以为频点顺序是「从 0 到最大频率」。DFT 的排列是
$0,1,\dots,N/2,-N/2+1,\dots,-1$，中间那个折返点最容易被忽略。
\end{know}

\begin{know}{H6}{radix-2：位反转与蝶形}{框架：math/fourier.hpp}
\field{是什么} 按频率下标的二进制位反转做一次重排，然后做 $\log_2N$ 层蝶形，每层步长 $2,4,\dots,N$。

\field{为什么需要} 把 $O(N^{2})$ 的求和拆成两半递归，是「分治」在变换上的原型；
它也是「算法结构决定常数」的样板。

\field{关键结论} 位反转的原因：递归拆偶/奇下标后，叶子顺序是「下标反转二进制」。
$N=1$ 合法（恒等）。成立条件：$N=2^{k}$；旋转因子必须\textbf{直接算 $\cos/\sin$ 建表}——
用递推 $w^{k+1}=w^{k}\cdot w$ 会把误差累到 $N\eps$（$N=2^{20}$ 时约 $2\times10^{-10}$，跨过 $10^{-12}$ 判据）。

\field{判据} 与 DFT 对拍（H1）；$\log$-$\log$ 斜率不是判据（这是精确算法，不是逼近）。

\field{前置} H1、A2。

\field{常见误解} 以为「FFT 是一种近似」。除了浮点舍入，FFT 与 DFT 给出\textbf{完全相同}的结果；
速度来自因子分解，不来自截断。
\end{know}

\begin{know}{H7}{复杂度：$O(N^{2})$ 与 $O(N\log N)$ 的量级差}{框架：math/fourier.hpp}
\field{是什么} DFT 需 $N^{2}$ 次复乘；FFT 需约 $\frac N2\log_2N$ 次蝶形（每次一个复乘两个复加）。

\field{为什么需要} 它解释「为什么 FFT 是信号处理的起点」：$N=2^{20}$ 时约 $10^{12}$ 对 $2\times10^{7}$，
差四个数量级。

\field{关键结论} 比值约 $2N/\log_2N$。成立条件：$N$ 为 2 的幂；常数因子未计。

\field{判据} 无（复杂度是量级陈述，不进 ctest）；实测只记一笔不设阈值。

\field{前置} H6。

\field{常见误解} 以为「$O(N\log N)$ 一定比 $O(N^{2})$ 快」。小 $N$（几十）时 DFT 的常数更小，
这也是库在小规模时用 DFT 作参照基线的原因。
\end{know}

\begin{know}{H8}{Parseval 与归一化因子}{框架：math/fourier.hpp}
\field{是什么} 在不归一化 + $1/N$ 逆变换的约定下：$\sum_n\abs{x_n}^{2}=\frac1N\sum_k\abs{X_k}^{2}$。

\field{为什么需要} 能量守恒是检验「变换有没有算错」的全局判据——它不依赖任何单个频点的精度。

\field{关键结论} 因子 $1/N$ 来自 H2 的逆变换归一化；若改用酉归一化（正逆各除 $\sqrt N$），因子变成 1。
成立条件：同一约定下两边比较；$\mathrm{abs2}$ 求和（不取模开方）。

\field{判据} 两边比值与 1 的差 $\le10^{-12}$。

\field{前置} H1、H2。

\field{常见误解} 以为 Parseval 是「近似成立的能量关系」。它是精确恒等式（离散、有限、无截断），
浮点上只差舍入。
\end{know}

\begin{know}{H9}{傅里叶级数：系数就是归一化谱}{框架：math/series.hpp}
\field{是什么} $c_k=\frac1N\sum_{n=0}^{N-1}x_ne^{-2\pi ikn/N}=X_k/N$，$k\in\Z$；
布局为 $(2K+1)\times1$，第 $r$ 行对应 $k=r-K$。

\field{为什么需要} 把「周期函数的频谱」与「离散 DFT」接上：级数系数不是新对象，就是 DFT 除 $N$。

\field{关键结论} 手算锚点：$x_n=\cos(2\pi\cdot3n/N)\Rightarrow c_{\pm3}=1/2$，其余 $0$。
截断重构 $f_K(t)=\sum_{k=-K}^{K}c_ke^{2\pi ikt}$ 在节点上等于 $N\cdot\mathrm{IFFT}(S)$——
\textbf{乘 $N$ 不是笔误}，因为 IFFT 自带 $1/N$（H2）。成立条件：$N=2^{k}$、$0\le K$、$2K+1\le N$；
截断窗口必须\textbf{对称}（单边截断得到的是解析信号，实部只剩一半）。

\field{判据} 漏乘 $N$ 的病态要能被幅度判据抓住（$f_K$ 与实际值差一个 $N$ 倍）。

\field{前置} H1、H2。

\field{常见误解} 以为 $c_k$ 与 $X_k$ 是「两个不同的量」。差别只有一个常数 $1/N$；
真正的差别在\textbf{索引约定}（$k$ 取负、居中存放）。
\end{know}

\begin{know}{H10}{截断与 Gibbs 过冲}{框架：math/series.hpp}
\field{是什么} 用前 $K$ 个谐波重构有跳变的周期函数时，跳变附近出现过冲，峰值趋于
$\frac2\pi\int_0^{\pi}\frac{\sin t}{t}dt=1.08948987\ldots$，即约 $8.95\%$ 的相对过冲。

\field{为什么需要} 它是「逐点收敛」与「一致收敛」分家的经典例子；也是「误差不随 $K$ 减小」的唯一常见例外。

\field{关键结论} 过冲幅度\textbf{与 $K$ 无关}（只与跳变高度有关），过冲宽度 $\sim1/K$——
所以 $L^{2}$ 误差趋于 0（H11），而 $\sup$ 误差不趋于 0。成立条件：函数有第一类间断。

\field{判据} 方波重构的峰值比 $\approx1.0895$，且 $K$ 增大时该比值不降；「不随 $K$ 减小」本身就是判据。

\field{前置} H9。

\field{常见误解} 以为「加更多项就能消除振铃」。可以把振铃\textbf{挤窄}，不能压平；
要压平必须换求和法（Fejér / Lanczos $\sigma$ 因子），那是改变函数而非提高分辨率。
\end{know}

\begin{know}{H11}{谱收敛：光滑度决定衰减速度}{框架：math/series.hpp}
\field{是什么} 若周期函数 $x$ 有 $p$ 阶连续导数（$x\in C^{p}$），则 $\abs{c_k}=O(\abs{k}^{-p-1})$。

\field{为什么需要} 它把「函数有多光滑」翻译成「系数掉得多快」，是选 $K$ 的唯一依据。

\field{关键结论} 方波 $p=0\Rightarrow\abs{c_k}=O(1/\abs{k})$，截断误差 $O(1/K)$；
解析函数则指数衰减。成立条件：$p$ 阶导数连续（分段光滑时按最差的那段算）。

\field{判据} 双对数坐标下 $\abs{c_k}$ 的斜率约 $-(p+1)$。

\field{前置} H9。

\field{常见误解} 以为「光滑函数的系数一定很快衰减」。$C^{1}$ 但二阶导数跳变的函数，
衰减只有 $O(k^{-3})$；指数衰减要求\textbf{解析}，比「无限可微」更强。
\end{know}

\section{概率：特征函数是「方向相反的那个傅里叶变换」}

\subsection*{串讲三：为什么 DFT 是 $-i$，特征函数是 $+i$}

这是全库唯一一处\textbf{故意反号}的地方（D20）。两边都在做「投影到 $e^{i\omega t}$」，
但一边是分析信号（把时间函数拆成频率成分），一边是刻画分布（把概率质量映射成
$\R$ 上的函数）。符号差异的具体后果很实在：

\begin{center}
\begin{tabular}{L{3.2cm}L{4.9cm}L{5.1cm}}
\toprule
 & DFT（$e^{-2\pi ikn/N}$） & 特征函数（$e^{+itx}$） \\
\midrule
反号后谁变 & $X_k\to X_{-k}$（谱镜像） & $\varphi'(0)$ 从 $+iE[X]$ 变 $-iE[X]$ \\
模长判据能否发现 & 能（谱位置反） & \textbf{不能}（共轭后模长相同） \\
修法 & 索引回卷 & 只能测\textbf{实/虚两槽} \\
\bottomrule
\end{tabular}
\end{center}

所以 P5 那条「接缝测试」必须用\textbf{复值概率质量函数}：实 pmf 下「只翻频轴」与「翻轴再共轭」
数值恰好相同，测试测不出错。

\begin{know}{P1}{特征函数与 $+i$ 约定}{框架：math/prob.hpp}
\field{是什么} $\varphi_X(t)=E[e^{itX}]=\int e^{itx}p(x)\,dx=\sum_k p_ke^{itx_k}$。

\field{为什么需要} 它把「独立和的分布」变成「函数相乘」（P3），把「矩」变成「导数」（P4）。

\field{关键结论} 与傅里叶变换方向\textbf{相反}（H1 用 $-i$）。成立条件：$E\abs{X}$ 存在时 $\varphi$
可微；$\varphi(0)=1$、$\abs{\varphi(t)}\le1$、$\varphi$ 一致连续。

\field{判据} 矩的判据必须分实/虚槽（P4），不能只看模。

\field{前置} C3、H1。

\field{常见误解} 以为「特征函数就是傅里叶变换」。数学上是，但\textbf{符号约定相反}；
照抄 FFT 的符号会得到 $\overline{\varphi}$，而模长判据完全看不出来。
\end{know}

\begin{know}{P2}{高斯的特征函数}{框架：math/prob.hpp}
\field{是什么} $X\sim N(\mu,\sigma^{2})\Rightarrow\varphi(t)=e^{i\mu t-\sigma^{2}t^{2}/2}$。

\field{为什么需要} 它是唯一「特征函数仍是高斯型」的分布，也是中心极限定理的解析支点。

\field{关键结论} $\mu=0$ 时 $\varphi$ 实、偶、非负；$\mu=0,\sigma=1$ 时
$\varphi(t)=e^{-t^{2}/2}$ \textbf{恰好等于}标准正态密度函数——高斯是傅里叶变换的不动点（差常数）。
成立条件：$\sigma>0$。

\field{判据} 高斯不动点：取 $dx\in\{0.5,0.4,0.3,0.25\}$ 四点，残差随 $dx$ 的衰减斜率与
$\exp(-\pi/dx^{2})$ 的预测一致。

\field{前置} P1。

\field{常见误解} 以为「不动点意味着 $\varphi=\varphi$ 恒成立」。它只在 $\mu=0,\sigma=1$ 且按
$\hat f(\xi)=\int f(x)e^{-2\pi i\xi x}dx$ 归一化时成立；换约定就多出常数与相位。
\end{know}

\begin{know}{P3}{卷积定理（概率版）}{框架：math/prob.hpp}
\field{是什么} $X\perp Y\Rightarrow\varphi_{X+Y}(t)=\varphi_X(t)\varphi_Y(t)$，等价于
$p_{X+Y}=p_X*p_Y$。

\field{为什么需要} 它是「独立随机变量求和」的唯一省力算法；也是 H3 在概率侧的同一条定理。

\field{关键结论} 成立条件：$X,Y$ 独立（不相关不够）；$\varphi$ 存在（即分布有定义）。

\field{判据} 与直接卷积对拍；离散情形用 H3 的循环卷积要注意补零（H4）。

\field{前置} P1、H3。

\field{常见误解} 以为「卷积定理说明两个和的分布等于两个密度相乘」。恰好\textbf{反了}：
是\textbf{特征函数}相乘、\textbf{密度}卷积。
\end{know}

\begin{know}{P4}{矩即导数：$i^{k}$ 因子与实虚分槽}{框架：math/prob.hpp}
\field{是什么} $\varphi^{(k)}(0)=i^{k}E[X^{k}]$。具体地：$\varphi'(0)=i\mu$（纯虚），
$\varphi''(0)=-(\mu^{2}+\sigma^{2})$（纯实）。

\field{为什么需要} 它把「求矩」变成「在 0 处求导」，是从特征函数反推分布的最短路径。

\field{关键结论} 奇阶导数带 $i^{\text{奇}}$ 因子，故落在虚部；偶阶落在实部。
成立条件：$E\abs{X}^{k}<\infty$；$0$ 处可导。

\field{判据} 判据必须分实/虚两槽比较；若整体用 $\mathrm{abs2}$，符号错误会被抹平。

\field{前置} P1。

\field{常见误解} 以为「$\varphi$ 在 0 处的斜率就是均值」。斜率是 $i\mu$，差一个 $i$；
把它当实数会得到 0（因为 $\mu$ 在虚部）。
\end{know}

\begin{know}{P5}{接缝：只翻频轴、不共轭}{框架：math/prob.hpp}
\field{是什么} 连续情形 $\varphi_X(2\pi\xi)=P(-\xi)$，其中 $P$ 是 $p$ 的傅里叶变换；
离散情形 $\varphi_X(2\pi k/N)=\mathrm{FFT}[p]_{(N-k)\bmod N}$。

\field{为什么需要} 库同时有 FFT（$-i$）与特征函数（$+i$），两者必须在某处对齐；
对齐点错了，整条「分布 $\leftrightarrow$ 频谱」的桥就歪了。

\field{关键结论} 正确做法是\textbf{只翻频轴}（索引回卷），\textbf{不额外共轭}。
成立条件：$p$ 为实 $pmf$ 时 $P(-\xi)=\overline{P(\xi)}$ 恰好成立，所以错误的实现看起来也对。

\field{判据} 必须用\textbf{复值} pmf 测：实 pmf 下「翻轴」与「翻轴+共轭」数值相同，测不出错。

\field{前置} P1、H1、H5。

\field{常见误解} 以为「反正实信号有共轭对称，共不共轭一个样」。共轭对称是\textbf{实信号}的性质，
一旦输入变成复数（比如已经乘过相位），这条性质消失，错误立刻显形。
\end{know}

\begin{know}{P6}{经验特征函数}{框架：math/prob.hpp}
\field{是什么} $\hat\varphi(t)=\frac1n\sum_{j=1}^{n}e^{itx_j}$，即把样本当经验分布的 pmf 代入 P1。

\field{为什么需要} 它是从数据估计特征函数的标准做法；大数律保证 $\hat\varphi\to\varphi_X$。

\field{关键结论} 误差量级 $O(1/\sqrt n)$（均值型统计量的标准差）；成立条件：样本独立同分布。
判据侧取 $3/\sqrt n$ 作容差。

\field{判据} $\abs{\hat\varphi(t)-\varphi(t)}\le3/\sqrt n$ 量级；同 seed 逐位可复现。

\field{前置} P1、P7。

\field{常见误解} 以为「样本量翻倍误差就减半」。$1/\sqrt n$ 意味着要误差减半需要\textbf{四倍}样本；
这个尺度差是很多「加数据没用」的根源。
\end{know}

\begin{know}{P7}{Box--Muller：从均匀到正态}{框架：math/prob.hpp}
\field{是什么} 取 $u_1,u_2\sim U(0,1]$ 独立，$r=\sqrt{-2\ln u_1}$，$\theta=2\pi u_2$，
则 $z_0=r\cos\theta$、$z_1=r\sin\theta$ 独立且 $\sim N(0,1)$。

\field{为什么需要} 它是「用可生成的均匀随机数造正态」的标准做法，也是「极坐标 + 逆变换」的经典组合。

\field{关键结论} 正确性：$R^{2}$ 服从参数 $1/2$ 的指数分布（即 $\chi^{2}_{2}$），$\Theta$ 均匀，
所以 $r$ 是「指数分布逆变换」；两个输出天然独立。成立条件：$u_1\in(0,1]$——
$u_1=0$ 会让 $\ln0=-\infty$ 污染均值。

\field{判据} $N=10^{6}$ 时 $\abs{\text{mean}}\le3/\sqrt N$、$\abs{\text{var}-1}\le3\sqrt{2/N}$；
同 seed 逐位可复现。

\field{前置} P1。

\field{常见误解} 以为「$r$ 该从 $N(0,1)$ 里取」。$r$ 是半径、服从指数分布；
把半径和角度都取成正态会得到一个错误的分布（而且看起来也像正态）。
\end{know}

\begin{know}{P8}{$\mathrm{erf}$、$\Phi$ 与尾部互补}{框架：math/special.hpp}
\field{是什么} $\mathrm{erf}(x)=\frac{2}{\sqrt\pi}\int_0^{x}e^{-t^{2}}dt$，
$\Phi(x)=\frac{1+\mathrm{erf}(x/\sqrt2)}{2}$，密度 $\varphi_{\text{pdf}}(x)=e^{-x^{2}/2}/\sqrt{2\pi}$。

\field{为什么需要} 正态的分布函数是「查表」的全部内容；尾部概率是显著性检验的语言。

\field{关键结论} 左尾必须走互补：$\Phi(-x)=\mathrm{erfc}(x/\sqrt2)/2\ (x>0)$——
因为 $x\gtrsim5.9$ 时 $\mathrm{erf}(x)=1-2.5\times10^{-16}$，只差 1 ulp，
用 $1-\mathrm{erf}$ 算左尾会得到 0。对称性 $\Phi(-x)=1-\Phi(x)$ 数学上成立但\textbf{数值上不能用}。
成立条件：$\sigma>0$ 且有限。

\field{判据} $\mathrm{erf}/\mathrm{erfc}$ 与语言内建差 $\le2$ ulp；$\abs{\Phi(\Phi^{-1}(u))-u}\le10^{-14}$
对 $u\in[10^{-12},1-10^{-12}]$。

\field{前置} A6、C5。

\field{常见误解} 以为「$\mathrm{erfc}=1-\mathrm{erf}$ 直接实现就行」。定义恒等，数值不恒等：
相消会把尾部信息全部抹掉（A6）。
\end{know}

\begin{know}{P9}{分位数：Acklam 近似 + 一步 Newton}{框架：math/special.hpp}
\field{是什么} 求 $z$ 使 $\Phi(z)=u$。先用 Acklam 有理逼近给初值 $z_0$（相对误差约 $1.15\times10^{-9}$），
再走一步 Newton：$z_1=z_0-\big(\Phi(z_0)-u\big)/\varphi_{\text{pdf}}(z_0)$。

\field{为什么需要} 分位数是「从概率反查阈值」的唯一工具（置信区间、显著性、采样）。

\field{关键结论} Newton 的收敛是二次的，前提是初值已落在收敛域内且 $\varphi(z_0)\ne0$（尾部会趋 0，
故 $u$ 必须严格在 $(0,1)$ 开区间）。成立条件：$u$ 不取端点（端点对应 $\pm\infty$，库不返回）。

\field{判据} $\abs{\Phi(z)-u}\le10^{-14}$，$u\in[10^{-12},1-10^{-12}]$；
锚点 $u=0.975\Rightarrow z\approx1.959963984540054$。

\field{前置} P8。

\field{常见误解} 以为「近似公式够了，不必再迭代」。Acklam 单独只有约 $10^{-9}$，
而 $\Phi$ 本身能到 ulp 级——不迭代就把上游精度白白丢掉三个数量级。
\end{know}

\begin{know}{P10}{尾部相消与判据的位置}{框架：math/special.hpp}
\field{是什么} 尾部（$|x|$ 大）概率极小，任何「大数减小数」的写法都会丢光有效位；
所以判据不能只写在「函数值」上，还要写在「反函数复合」上。

\field{为什么需要} 尾部是「精度最容易假通过」的地方：中心区域对拍全绿，尾部错得离谱。

\field{关键结论} 两条互补判据：$\mathrm{erfc}$ 与内建对拍（绝对差 $\le2$ ulp）覆盖尾部；
$\Phi(\Phi^{-1}(u))$ 复合判据覆盖「往返一致性」。成立条件：$u$ 在 $[10^{-12},1-10^{-12}]$ 内，
再往外 $\Phi^{-1}$ 的绝对误差会让复合判据失真。

\field{判据} 复合误差 $\le10^{-14}$；$\mathrm{erfc}(+\infty)=0$、$\mathrm{erfc}(-\infty)=2$、$\mathrm{erf}(\pm\infty)=\pm1$。

\field{前置} P8、P9、A6。

\field{常见误解} 以为「函数值与反函数值都测了就万事大吉」。两者可能\textbf{一致地错}
（同一个相消 bug 让两边偏移相同），必须再加一条与外部实现的绝对对拍。
\end{know}

\section{微分：从 $\eps^{2}=0$ 到伴随传播}

\subsection*{串讲四：为什么「前向一次得列、反向一次得行」}

这件事看着像工程技巧，其实全是维度：

\begin{itemize}[leftmargin=1.4em]
  \item 前向：把 $\eps$ 塞进第 $j$ 个输入分量 $x_j$。仿射部分立刻吐出
        $\partial f/\partial x_j$——\textbf{一次一个输入}，是 $J$ 的一列。
  \item 反向：从标量 $L$ 出发，把「上游的伴随」按链式法则往回乘——\textbf{一次一个输出}，是 $J$ 的一行。
\end{itemize}

于是代价对偶：$n$ 个输入、$m$ 个输出时，前向 $O(n)$ 次、反向 $O(m)$ 次。
深度学习里 $m=1$（损失是标量）而 $n$ 是百万参数，所以必须反向。这不是「反向模式更先进」，
只是 $m\ll n$ 时它更便宜。

\begin{know}{D1}{对偶数：$\eps^{2}=0$ 的代数}{已实现：autodiff/dual.hpp}
\field{是什么} $z=v+d\eps$，规定 $\eps^{2}=0$ 且 $\eps\ne0$。注意 $\eps$ 不是 $i$：
$i^{2}=-1$，$\eps^{2}=0$。

\field{为什么需要} 它把「求导」变成「在同一个数域里做算术」——不需要符号系统，也不需要差分步长。

\field{关键结论} 对任意光滑 $f$，$f(v+d\eps)=f(v)+d\,f'(v)\eps$：因为 $h^{2}=0$，
\textbf{高阶项精确消失}，所以链式法则成为代数恒等式（不是极限意义下的近似）。成立条件：
$f$ 在 $v$ 的邻域内解析（有 Taylor 展开）；$T$ 必须是\textbf{交换环}——四元数不满足，不得进此域。

\field{判据} 二阶导数用嵌套 $\mathrm{Dual\langle Dual\rangle}$：$v+d_1\eps_1+e\eps_2+(\partial^{2})\eps_1\eps_2$。

\field{前置} 无。

\field{常见误解} 以为「$\eps$ 是个很小的数」。它不是数，是一个幂零生成元；
写成 $10^{-8}$ 就退化成数值差分，误差从机器精度掉到 $O(\sqrt\eps)$。
\end{know}

\begin{know}{D2}{传播规则与链式法则}{已实现：autodiff/dual.hpp}
\field{是什么} $(a+b\eps)(c+e\eps)=ac+(ae+bc)\eps$；
$(a+b\eps)/(c+e\eps)=\frac ac+\frac{bc-ae}{c^{2}}\eps$（$c\ne0$）。
原函数：$\sin:( \sin a,\ \cos a\cdot b)$、$\exp:(\exp a,\ \exp a\cdot b)$、$\ln:(\ln a,\ b/a)$、
$\sqrt:(\sqrt a,\ b/(2\sqrt a))$。

\field{为什么需要} 只要每种原函数都满足「值 + 导数」二元传播，整条表达式树就自动正确——
这是把微积分降级成算术的关键。

\field{关键结论} 连乘的导数用前后缀积可做到 $O(n)$ 而避免零因子问题：
$dY=\sum_i\big(\prod_{j\ne i}v_j\big)dx_i$。成立条件：$c\ne0$（除零由 IEEE 传播，标量层不判，D12）。

\field{判据} 与解析导数对拍；$n$ 元连乘的 $O(n)$ 版本与朴素 $O(n^{2})$ 版本互相对拍。

\field{前置} D1。

\field{常见误解} 以为「除法要单独实现一套导数」。它完全由乘法与可逆元推出；
手写公式容易在 $c$ 接近 0 时与乘法版本不一致。
\end{know}

\begin{know}{D3}{雅可比布局：行是输出、列是输入}{已实现：autodiff/jacobian.hpp}
\field{是什么} 分子布局：$J\in\R^{m\times n}$，$J_{ij}=\partial y_i/\partial x_j$（D19 第二次修订）；
链式法则取自然序 $J_{f\circ g}=J_f\cdot J_g$。

\field{为什么需要} 布局不统一，全库的「转置」会到处漏；链式法则的方向一错，反传全错。

\field{关键结论} 标量输出时 $J$ 是 $1\times n$ \textbf{行}向量，而梯度 $\nabla f=J^{\mathsf T}$ 是 $n\times 1$ 列——
两者\textbf{不是同一个东西}，梯度与被微对象同形。成立条件：向量一律 $n\times1$（D19）。

\field{判据} 差分校验 $J_{ij}\approx\frac{f_i(x+he_j)-f_i(x-he_j)}{2h}$，$h=10^{-6}$ 时误差约 $10^{-10}$，
判据取差 $<10^{-8}$。

\field{前置} L1、L2。

\field{常见误解} 以为「分母布局只是写法的偏好」。两种布局互为转置，链式法则的乘法顺序相反
（$J_gJ_f$）；混用一次就会得到形状对、数值错的结果。
\end{know}

\begin{know}{D4}{种子 sweep：一次填一列}{已实现：autodiff/jacobian.hpp}
\field{是什么} 令 $x=x_0+e_j\eps$，则 $f(x)=f(x_0)+(\partial f/\partial x_j)\eps$；
扫遍 $j=1..n$ 即得完整 $J$，共 $n$ 次前向。

\field{为什么需要} 它是前向模式的实际算法，也是「代价 = 输入维数」这个说法的出处。

\field{关键结论} 第 $j$ 次 sweep 只填第 $j$ 列；形状校验必须\textbf{逐维}（只比一个维度会退化成恒真）。
成立条件：$x$ 为 $n\times1$、输出 $m\times1$。

\field{判据} 与 D3 的差分校验对拍；对角/逐元素函数的 $J$ 应为 $\mathrm{diag}(\varphi')$，
广播（$f$ 为 $1\times1$）时 $J=\varphi'(a)\cdot\mathbf{1}$。

\field{前置} D1、D3。

\field{常见误解} 以为「一次 sweep 能得整行」。前向一次只给一列；要整行必须反向（D5）。
方向搞反会让整体代价从 $O(n)$ 变成 $O(n^{2})$ 而不报错。
\end{know}

\begin{know}{D5}{反向模式与伴随传播}{已实现：autodiff/grad\_graph.hpp}
\field{是什么} 设 $v=\partial L/\partial z$ 为「上游伴随」，则
$\partial L/\partial n=\sum_{m:n\to m}(\partial L/\partial m)(\partial m/\partial n)$。

\field{为什么需要} $m=1$ 时一次反向得到整个梯度，代价与参数规模成线性——这是训练可行性的全部基础。

\field{关键结论} 三条要点：(1)fan-out 处必须\textbf{累加}（同一节点被多条路径用到）；
(2)局部导数按算子给：MatMul 是 $\partial L/\partial W=vx^{\mathsf T}$、$\partial L/\partial x=W^{\mathsf T}v$；
(3)形状唯一确定公式（$vx^{\mathsf T}$ 必为 $m\times n$）。成立条件：图是有向无环；无 tape 时 fan-out 会重算导致指数爆炸。

\field{判据} 恒等图 $L=x$ 的梯度必须恰为 $1$；差分校验 $h=10^{-6}$、误差约 $1\times10^{-10}$。

\field{前置} D3、L10。

\field{常见误解} 以为「反向就是前向的逆运算」。它不是逆矩阵，是\textbf{转置}（伴随）传播：
每一步乘的是局部雅可比的转置，而不是它的逆。
\end{know}

\begin{know}{D6}{Lipschitz 条件与解的存在唯一}{骨架：math/ode.hpp}
\field{是什么} 初值问题 $y'=f(t,y)$、$y(t_0)=y_0$；若 $f$ 对 $y$ 满足 Lipschitz 条件
$\abs{f(t,y_1)-f(t,y_2)}\le L\abs{y_1-y_2}$，则局部解存在且唯一（Picard--Lindelöf）。

\field{为什么需要} 没有唯一性，任何数值解都失去「在对什么东西逼近」的前提。

\field{关键结论} Lipschitz 弱于连续可微，强于连续；只有连续性（Peano）时解可能不唯一。
契约要求 $f$ 光滑；\textbf{跳变不处理}（需要事件检测，不做）。

\field{判据} 无（骨架未实现）。

\field{前置} A5。

\field{常见误解} 以为「右端连续就一定有唯一解」。经典反例 $y'=\sqrt{\abs y},y(0)=0$ 有无穷多解；
差的就是 Lipschitz。
\end{know}

\begin{know}{D7}{Euler 与 RK4 的阶}{骨架：math/ode.hpp}
\field{是什么} Euler：$y_{n+1}=y_n+hf(t_n,y_n)$，1 阶。RK4：
$k_1=f(t,y)$，$k_2=f(t+h/2,y+\frac h2k_1)$，$k_3=f(t+h/2,y+\frac h2k_2)$，$k_4=f(t+h,y+hk_3)$，
$y_{n+1}=y+h(k_1+2k_2+2k_3+k_4)/6$，4 阶。

\field{为什么需要} 它们是「局部误差阶 $\to$ 全局误差阶」的标准例子，也是收敛阶判据（A5）的常用靶子。

\field{关键结论} Euler 局部 $O(h^{2})$、全局 $O(h)$（$h$ 减半误差降 2 倍）；
RK4 局部 $O(h^{5})$、全局 $O(h^{4})$（$h$ 减半误差降 16 倍），每步 4 次 $f$。
成立条件：$f$ 足够光滑且步长进入渐近区。

\field{判据} 三次 $h$ 减半的误差比值（2 与 16），落区间内才算过。

\field{前置} D6、A5。

\field{常见误解} 以为「RK4 一定比 Euler 准」。$h$ 很大时 RK4 的 4 次求值可能落在解的外面，
反而发散；阶是\textbf{渐近}性质。
\end{know}

\begin{know}{D8}{辛结构与 Verlet 的能量行为}{骨架：math/ode.hpp}
\field{是什么} 力学系统写成 $\dot q=v$、$\dot v=a(q)$。辛 Euler：$v_{n+1}=v_n+ha(q_n)$、
$q_{n+1}=q_n+hv_{n+1}$（用\textbf{新}速度）。Verlet：半步速度—整步位置—半步速度
（$v_{1/2}=v_n+\frac h2a(q_n)$、$q_{n+1}=q_n+hv_{1/2}$、$v_{n+1}=v_{1/2}+\frac h2a(q_{n+1})$）。

\field{为什么需要} 长时间积分时，普通高阶方法会让能量单调漂移；保结构方法让能量\textbf{有界振荡}。

\field{关键结论} 辛方法精确保持的是\textbf{扰动哈密顿量} $H+O(h^{2})$，故能量误差不发散。
Verlet 2 阶、每步 2 次加速度；辛 Euler 1 阶但已有界。成立条件：$H=\frac12\abs v^{2}+U(q)$，
力由 $a=-\nabla U$ 给出（保守系统）。

\field{判据} 判能量\textbf{包络}（上下界是否随时间增长），不判单步误差；
比较辛 Euler 与 RK4 的长时间能量曲线是标准用例。

\field{前置} D7。

\field{常见误解} 以为「阶更高就更守恒」。RK4 阶高但不保辛，长时间能量会\textbf{系统性漂移}；
这是「结构 > 阶数」最著名的反例。
\end{know}

\section{几何与李群：同一个旋转的三件衣服}

\subsection*{串讲五：$\mathrm{hat}$、四元数、$SO(3)$ 为什么是一件事}

三种表示都在描述「绕某轴的转动」，差别只在\textbf{存什么、怎么复合}：

\begin{center}
\begin{tabular}{L{2.3cm}L{3.2cm}L{3.4cm}L{4.2cm}}
\toprule
表示 & 存什么 & 复合 & 代价 / 陷阱 \\
\midrule
旋转矩阵 $R\in SO(3)$ & 9 个数 & 矩阵乘 & 无奇点，但 9 个数里只有 3 个自由度；正交性会漂 \\
李代数 $\hat u\,\theta\in\mathfrak{so}(3)$ & 3 个数 & 需要 $\exp$ 才能回到群 & 最省，但 $\theta$ 有 $2\pi$ 周期、$\pi$ 附近反解不稳 \\
单位四元数 $q\in S^{3}$ & 4 个数 & 四元数乘 & 无双覆盖歧义但有 $\pm q$；插值最自然 \\
\bottomrule
\end{tabular}
\end{center}

桥梁是两条：$\exp(\theta\hat u)$ 给出 Rodrigues（G2），而四元数的 $q=\cos\frac\theta2+\sin\frac\theta2\hat u$
与它满足 $R(q)=\exp(\theta\hat u)$（G6）。\textbf{半角}不是技巧，是双覆盖的必然：
$q$ 与 $-q$ 必须映到同一个 $R$，所以角度只能吃一半。

\begin{know}{G1}{叉积、hat 与反对称矩阵}{已实现：linalg/so3.hpp}
\field{是什么} $a\times b=(a_yb_z-a_zb_y,\ a_zb_x-a_xb_z,\ a_xb_y-a_yb_x)$；
$\mathrm{Hat}(u)$ 是满足 $\mathrm{Hat}(u)b=u\times b$ 的 $3\times3$ 反对称矩阵，
即 $\mathrm{Hat}(u)=\begin{bmatrix}0&-u_3&u_2\\u_3&0&-u_1\\-u_2&u_1&0\end{bmatrix}$。

\field{为什么需要} 它把「叉积」变成「矩阵作用」，于是旋转可以写成指数——这就是李代数到李群的入口。

\field{关键结论} 性质：$\mathrm{Hat}(u)^{\mathsf T}=-\mathrm{Hat}(u)$、
$\mathrm{Hat}(\lambda u)=\lambda\,\mathrm{Hat}(u)$（\textbf{线性，绝不归一化}）、
$\mathrm{Hat}(u)^{2}=uu^{\mathsf T}-\norm{u}^{2}I$、$\abs{a\times b}=\abs{a}\abs{b}\sin\angle$。
成立条件：入参出参一律 $3\times1$ 列向量（D19），形状不对抛错而不是静默接受。

\field{判据} $\mathrm{Hat}(u)b$ 与显式叉积逐分量对拍；$\mathrm{Hat}$ 的反对称性逐元素查。

\field{前置} L10、L3。

\field{常见误解} 以为「$\mathrm{Hat}$ 该把 $u$ 归一化」。归一化会把 $\theta$ 从矩阵里抹掉，
$\exp$ 之后只能得到单位角度；线性性是整套记号成立的前提。
\end{know}

\begin{know}{G2}{立方恒等式与 Rodrigues 公式}{已实现：linalg/so3.hpp}
\field{是什么} $\mathrm{Hat}(u)^{3}=-\norm{u}^{2}\mathrm{Hat}(u)$；于是
$\exp(\theta\,\mathrm{Hat}(\hat u))=I+\sin\theta\,\mathrm{Hat}(\hat u)+(1-\cos\theta)\,\mathrm{Hat}(\hat u)^{2}$。

\field{为什么需要} 它是 $SO(3)$ 的闭式指数映射——不需要矩阵级数，也不需要特征分解。

\field{关键结论} 立方恒等式让级数在三次之后闭合，故 $\exp$ 只有三项。单位轴时化为
$\mathrm{Hat}(\hat u)^{3}=-\mathrm{Hat}(\hat u)$。成立条件：$\norm{\hat u}=1$（否则把 $\theta$ 换成 $\theta\norm u$
并把轴换成 $\hat u=u/\norm u$）。源码实测残差 $2.2\times10^{-16}$。

\field{判据} $\norm{\exp(\theta\hat u)-R_{\text{显式}}}$；立方恒等式残差；
但\textbf{整体符号翻转抓不到}（$\hat u\to-\hat u,\theta\to-\theta$ 给同一 $R$）。

\field{前置} G1、C3。

\field{常见误解} 以为「$R$ 唯一决定轴和角」。$R$ 只决定 $\{\pm\hat u,\ \theta\leftrightarrow2\pi-\theta\}$ 这一对等价类；
反解必须额外定约定（通常取 $\theta\in[0,\pi]$）。
\end{know}

\begin{know}{G3}{vee 与对数映射：用 atan2 而不是 arccos}{已实现：linalg/so3.hpp}
\field{是什么} $\mathrm{Vee}(K)=(K_{21},K_{02},K_{10})$，满足 $\mathrm{Vee}(\mathrm{Hat}(u))=u$；
反解
\[
\theta=\mathrm{atan2}\Bigl(\tfrac12\norm{\mathrm{Vee}(R-R^{\mathsf T})},\ \tfrac12(\mathrm{tr}R-1)\Bigr).
\]

\field{为什么需要} 把 $R$ 变回「轴 $\times$ 角」，才能在两个旋转之间插值、做误差度量。

\field{关键结论} 用 $\mathrm{atan2}$ 而不是 $\arccos$：后者在 $\theta\to0$ 时相对误差约 $\sqrt{\eps}/\theta$，
$\theta=10^{-8}$ 直接返回 0。成立条件：$\theta\ne\pi$（此时 $R-R^{\mathsf T}=0$，轴不可辨，
需要改用对称部分 $\frac{R+I}{2}$ 求轴）。

\field{判据} 往返 $\mathrm{Exp}(\mathrm{Log}(R))$ 与 $R$ 的差；$\theta$ 很小时（$10^{-8}$）仍应给出非零轴角。

\field{前置} G1、G2、A6。

\field{常见误解} 以为「$\arccos$ 与 $\mathrm{atan2}$ 只是写法不同」。在 $\theta$ 接近 0 或 $\pi$ 的两端，
两者的\textbf{有效位}差几个数量级；这是「同一个函数、两种数值行为」的典型。
\end{know}

\begin{know}{G4}{李括号：$[\mathrm{Hat}(a),\mathrm{Hat}(b)]=\mathrm{Hat}(a\times b)$}{已实现：linalg/so3.hpp}
\field{是什么} 矩阵换位子 $[A,B]=AB-BA$；对 $\mathfrak{so}(3)$ 有
$[\mathrm{Hat}(a),\mathrm{Hat}(b)]=\mathrm{Hat}(a\times b)$。

\field{为什么需要} 它说明「$\mathfrak{so}(3)$ 对换位子封闭」，即它确实是李代数；
也是「旋转不可交换」的代数表达。

\field{关键结论} 因子是 $1$ 不是 $2$（不同书的约定会差常数，本库取 1）。
成立条件：$a,b$ 为 $3\times1$。

\field{判据} 两边逐分量对拍（精确恒等式，容差到 $10^{-14}$ 即可）。

\field{前置} G1。

\field{常见误解} 以为「换位子小就说明两个旋转接近可交换」。换位子为零才是可交换（共轴）；
它的大小与角度无关地给出「转轴夹角」，不是「角度差」。
\end{know}

\begin{know}{G5}{四元数代数：唯一非交换来源是叉积}{已实现：quat/quaternion.hpp}
\field{是什么} $i^{2}=j^{2}=k^{2}=ijk=-1$；
$(w_1+\mathbf v_1)(w_2+\mathbf v_2)=(w_1w_2-\mathbf v_1\cdot\mathbf v_2)+
(w_1\mathbf v_2+w_2\mathbf v_1+\mathbf v_1\times\mathbf v_2)$。

\field{为什么需要} 四元数是最省的四维代数表示旋转：4 个数、无奇点、复合只需 16 次乘法。

\field{关键结论} 最后一项 $\mathbf v_1\times\mathbf v_2$ 是\textbf{唯一的非交换来源}；
去掉它就是复数。成立条件：$T$ 为浮点；四元数是除环（非零元可逆）但\textbf{非交换}，
所以不得进入自动微分的交换环域（D1）。

\field{判据} 结合律与单位元；$qq^{-1}=1$；$\abs{q_1q_2}=\abs{q_1}\abs{q_2}$（范数乘法性）。

\field{前置} G1、C1。

\field{常见误解} 以为「四元数乘法就是四个分量各自乘」。它含叉积项，顺序反了得到共轭旋转；
把四元数当向量逐分量算会「看起来像」但结果不是旋转。
\end{know}

\begin{know}{G6}{双覆盖与半角}{已实现：quat/quaternion.hpp}
\field{是什么} $q=\cos\frac\theta2+\sin\frac\theta2\hat u$ 对应旋转 $\theta$；
$q$ 与 $-q$ 给出同一个旋转矩阵，即映射 $S^{3}\to SO(3)$ 是\textbf{二对一}（双覆盖）。

\field{为什么需要} 它解释了半角、$\theta+2\pi$ 的等价、以及为什么四元数的「角度」只能取 $[0,\pi]$ 的一半。

\field{关键结论} $M(q)=M(-q)$；旋转复合对应四元数乘（同向）；共轭对应逆旋转。
成立条件：$\norm{q}=1$（单位四元数）才有纯旋转；非单位会带缩放。

\field{判据} $M(q)=M(-q)$ 逐元素；$\abs q=1$ 时 $M(q)$ 正交且 $\det=1$。

\field{前置} G2、G5。

\field{常见误解} 以为「四元数转动 $2\pi$ 就回到原状态」。四元数要 $4\pi$ 才回到自己
（$q\to-q$），这正是「自旋 $1/2$」的几何来源；用它插值时选短弧就是在处理这个歧义（G8）。
\end{know}

\begin{know}{G7}{四元数的逆与零判据}{已实现：quat/quaternion.hpp}
\field{是什么} $q^{-1}=q^{*}/\abs{q}^{2}$，$q^{*}=(w,-\mathbf v)$；单位四元数下 $q^{-1}=q^{*}$。

\field{为什么需要} 逆旋转、误差度量（$q_1^{-1}q_2$ 的角）、以及插值的中间量都用到它。

\field{关键结论} 判零用\textbf{四个分量精确全为 0}，不带容差——与矩阵求逆的容差判据（N9）不同。
成立条件：非零四元数。

\field{判据} $qq^{-1}=1$（到 $10^{-15}$）；零四元数抛异常而不是返回 inf。

\field{前置} G5。

\field{常见误解} 以为「四元数求逆和矩阵求逆一样该带容差」。矩阵的「接近奇异」有实际意义
（条件数大），四元数的「接近零」只是一个点；加容差会把合法的极小四元数误杀。
\end{know}

\begin{know}{G8}{球面插值：slerp 与 nlerp}{已实现：quat/quaternion.hpp}
\field{是什么} $\mathrm{slerp}(t)=\dfrac{\sin((1-t)\Omega)}{\sin\Omega}q_0+\dfrac{\sin(t\Omega)}{\sin\Omega}q_1$，
$\Omega=\arccos\abs{q_0\cdot q_1}$；退化时用线性插值再归一化（nlerp）。

\field{为什么需要} 关键帧动画、姿态平滑都用它；「匀速转动」是它唯一的正确性标准。

\field{关键结论} 短弧判据：$q_0\cdot q_1<0$ 时把 $q_1$ 翻号（否则绕远路）。
$\Omega$ 很小时 $\sin\Omega$ 出现在分母会放大误差，故 $d>0.9995$（$\Omega\ge0.0316$）时切 nlerp。
成立条件：$q_0,q_1$ 均已归一。

\field{判据} 插值后 $\abs q=1$；角速度恒定的用例（等步长 $t$ 给出的角增量应相等）。

\field{前置} G6、G7。

\field{常见误解} 以为「nlerp 就是 slerp 的近似、差别很小」。线性插值会让 $\abs q$ 掉到 0.857 附近，
再归一化后\textbf{角速度不均匀}——慢—快—慢，动画上看得出顿挫。
\end{know}

\begin{know}{G9}{从旋转矩阵反解四元数（Shepperd）}{已实现：quat/quaternion.hpp}
\field{是什么} 四个候选量形如 $\frac{1+R_{ii}-R_{jj}-R_{kk}}{4}$，取\textbf{最大}的那个开根作「大头」，
再由它推出其余三个分量。

\field{为什么需要} 从 $R$ 回到 $q$ 是刚体状态互换的必经步骤；单分支写法在 $\theta\to\pi$ 处精度崩。

\field{关键结论} 之所以要分支：当某个分量接近 0 时，从它反解会把误差放大到 $\sqrt{\eps}$ 量级；
选最大者保证分母不为小量。成立条件：$R$ 是正交且 $\det=1$（不是的话先做正交化）。

\field{判据} 往返 $M(\mathrm{Quat}(R))=R$；$\theta=\pi$ 附近（$\mathrm{tr}R\to-1$）单独加用例。

\field{前置} G6、G2。

\field{常见误解} 以为「用 $\sqrt{1+\mathrm{tr}R}/2$ 一个公式就够了」。$\mathrm{tr}R\to-1$ 时它就是 $0/0$；
所有教程里的「一个公式」都隐含了「$w$ 不小」的假设。
\end{know}

\begin{know}{G10}{角度公式：为什么是 $2\,\mathrm{atan2}(\abs{\mathbf v},w)$}{已实现：quat/quaternion.hpp}
\field{是什么} $\theta=2\,\mathrm{atan2}(\norm{\mathbf v},w)$（先令 $w\ge0$）；
两旋转之差 $\Omega=2\arccos\abs{q_0\cdot q_1}$。

\field{为什么需要} 角度是所有「误差/容差」的度量单位：姿态误差多大算收敛，全靠它。

\field{关键结论} 不用 $2\arccos w$：后者的相对误差约 $\sqrt{\eps}/\theta$，$\theta=10^{-8}$ 时返回 0。
成立条件：$q$ 已归一（否则 $\norm{\mathbf v}$ 与 $w$ 不在同一圆上）。

\field{判据} 小角度用例（$\theta=10^{-8}$）必须给出非零值；与 G3 的矩阵路径对拍。

\field{前置} G3、G6。

\field{常见误解} 以为「$\arccos$ 更直接所以更好」。在 0 附近 $\cos\theta\approx1-\theta^{2}/2$，
$\arccos$ 要开方两次才能分辨；$\mathrm{atan2}$ 用的是比值，天然避开这个相消。
\end{know}

\section{离散化与变分：连续方程怎么变成 $Au=b$}

这一节全部是 \textbf{未建}（只在 \texttt{NOW.md} 规划）。它的价值不在「实现细节」，
而在\textbf{把一道 PDE 拆成五步数学动作}——面试里能说清这五步的人远少于能背公式的人。

\begin{center}
\begin{tabular}{L{1.1cm}L{3.2cm}L{4.0cm}L{5.2cm}}
\toprule
步 & 数学动作 & 产物 & 为什么不能省 \\
\midrule
1 & 选网格、定自由度 & 节点集 $\{x_i\}$、编号表 & 编号决定矩阵的稀疏结构与带宽 \\
2 & 弱形式 + 分部积分 & $\int\nabla u\cdot\nabla v=\int fv$ & 降阶：二阶导挪到测试函数上 \\
3 & 单元刚度装配 & $K_e\to A$ & 每个单元的积分只算一次 \\
4 & 载荷 $b_i=\int f\varphi_i$ & $b$ & 用闭式积分才不引入额外误差 \\
5 & 边界条件 & 消元后的 $A,b$ & 不改边界，矩阵奇异（常值在零空间）\\
\bottomrule
\end{tabular}
\end{center}

\begin{know}{F1}{结构化网格与自由度编号}{未建：math/grid.hpp}
\field{是什么} 1D：$[a,b]$ 上 $n$ 个节点、$n-1$ 个单元，$h=(b-a)/(n-1)$。
2D：$n_x\times n_y$ 个节点、$(n_x-1)(n_y-1)$ 个矩形单元，每格切两个三角形，
节点按行主序编号 $p=j\,n_x+i$。

\field{为什么需要} 编号不是记账，是\textbf{矩阵结构}：行主序编号给出带宽 $O(n_x)$ 的带状矩阵，
换成对角编号带宽变 $O(n_x+n_y)$——同一个 PDE、不同的解算代价。

\field{关键结论} 自由度 = 节点（P1 元）；边界自由度是四条边上的节点。成立条件：矩形域、等距、固定对角线方向。

\field{判据} 手算表对拍：1D $n=5,[0,1]\Rightarrow$ 节点 $\{0,.25,.5,.75,1\}$、边界 $\{0,4\}$；
2D $4\times3\Rightarrow$ 12 节点、内点 2、边界 10、三角形 $2\cdot3\cdot2=12$。

\field{前置} L1。

\field{常见误解} 以为「单元数等于节点数」。1D 单元是 $n-1$、2D 单元是 $2(n_x-1)(n_y-1)$；
把单元当节点会导致刚度矩阵维度错。
\end{know}

\begin{know}{F2}{P1 单元刚度矩阵}{未建：math/assemble.hpp}
\field{是什么} 1D：$K_e=\frac1h\begin{bmatrix}1&-1\\-1&1\end{bmatrix}$；
2D 直角三角形：$K_e=\frac12\begin{bmatrix}2&-1&-1\\-1&1&0\\-1&0&1\end{bmatrix}$。

\field{为什么需要} 它是「$-\Delta$ 在分片线性空间上的矩阵表示」，整个有限元就是从它装配出来的。

\field{关键结论} 两个必查性质：\textbf{行和为 0}（常函数在核里，对应「常数不影响梯度」）、
\textbf{对称}（因为双线性形式对称）。成立条件：P1（分片线性）基函数；单元几何规则。

\field{判据} 1D $n=3,h=0.5$：$A=\frac1h\begin{bmatrix}1&-1&0\\-1&2&-1\\0&-1&1\end{bmatrix}$ 逐项对拍；
$A(i,j)=A(j,i)$ 要求浮点\textbf{精确}相等（同一条装配路径算两次）。

\field{前置} F1。

\field{常见误解} 以为「2D 的 $K_e$ 是 1D 的张量积」。直角三角形的 $K_e$ 不是矩形元的张量积形式；
用错会得到与解析解不符的收敛阶。
\end{know}

\begin{know}{F3}{弱形式与分部积分}{未建：math/assemble.hpp}
\field{是什么} 对 $-\Delta u=f$（Dirichlet 0），两边乘测试函数 $v$ 并积分，用分部积分：
\[
\int_\Omega\nabla u\cdot\nabla v\,dx=\int_\Omega fv\,dx\quad\forall v\in H^{1}_{0}.
\]

\field{为什么需要} 强形式要求 $u$ 二阶可导；弱形式只要求一阶可积——这才让分片线性函数成为合法解。

\field{关键结论} 分部积分把「一个二阶导」换成「两个一阶导的乘积」，同时\textbf{把边界项丢给自然边界条件}。
Dirichlet 条件是本质边界（要强加），Neumann 是自然边界（出现在右端）。
成立条件：$u,v\in H^{1}$；边界足够正则（分段光滑即可）。

\field{判据} 无（这一步是推导，不是计算）；但「行和为 0」（F2）就是弱形式在常数上的推论。

\field{前置} F2。

\field{常见误解} 以为「弱形式是强形式的近似」。两者对光滑解\textbf{完全等价}；弱形式只是把要求放宽到
更差的函数类，从而让离散解有定义。
\end{know}

\begin{know}{F4}{一致载荷与闭式积分}{未建：math/assemble.hpp}
\field{是什么} $b_i=\int f\varphi_i$。1D 线性 $f$：单元上 $\frac h6[2f_i+f_{i+1},\ f_i+2f_{i+1}]$；
2D 三角形用 $\int\lambda_1^{a}\lambda_2^{b}\lambda_3^{c}\,dA=2A\frac{a!\,b!\,c!}{(a+b+c+2)!}$。

\field{为什么需要} 载荷的误差会\textbf{直接进入解}；用「节点值代替积分」（集中质量）会把插值误差与
积分误差混在一起，让收敛阶无法归因。

\field{关键结论} 只接受次数 $\le2$ 的多项式 $f$（有闭式），其余需要求积器（不做）。
成立条件：$f$ 多项式；单元几何规则。

\field{判据} 1D 节点精确性（F6）：$u=x(1-x)$、$f=2$ 时 P1 解在节点上与解析解\textbf{逐点相等}（机器精度）——
载荷、边界、单元刚度任一错，这条立刻红。

\field{前置} F2、F3。

\field{常见误解} 以为「一致性质量矩阵不如集中质量好」。集中质量是对角矩阵、求解快，
但它\textbf{对应另一个离散问题}；混用会让「二阶收敛」的归因失效。
\end{know}

\begin{know}{F5}{Dirichlet 消元与保持对称}{未建：math/assemble.hpp}
\field{是什么} 对边界自由度：把对应行列清零、对角置 1、右端置为边界值，并把边界值对内部的贡献从 $b$ 里扣掉。

\field{为什么需要} 不处理边界，$A$ 有非平凡零空间（常向量），CG 必失败（F9）；
处理不当，$A$ 变不对称，CG 的整个理论前提崩掉。

\field{关键结论} 「行、列都清零 + 对角置 1」保持\textbf{对称}；只清行会破坏对称性。
成立条件：Dirichlet 值已知；边界自由度集合正确（F1）。

\field{判据} 处理后 $A$ 仍对称、对角非零、$x$ 的边界分量等于给定值；未加边界时内部节点行和为 0。

\field{前置} F1、F2。

\field{常见误解} 以为「把边界节点从矩阵里删掉」更干净。删行删列会改变未知量编号（要重排索引），
工程上更容易出错；清零法保留维度、保持结构。
\end{know}

\begin{know}{F6}{收敛阶 2 与 1D 的节点精确性}{未建：math/assemble.hpp}
\field{是什么} P1 元在 $L^{2}$ 下有二阶收敛：$\norm{u-u_h}_{L^{2}}\le Ch^{2}\abs{u}_{H^{2}}$。
1D 的特殊性：节点上的值\textbf{精确}（不用减 $O(h^{2})$）。

\field{为什么需要} 它是「离散解是否可信」的唯一客观判据；1D 的节点精确性是整条链最锋利的验收工具。

\field{关键结论} 阶 2 的来源：$L^{2}$ 误差 $\sim h\cdot$（插值误差 $O(h)$ 的一阶矩），
两个 $h$ 相乘。1D 节点精确的原因：Green 函数对称性使得节点误差的主项恒为零——这是 1D 的偶然，\textbf{2D 不成立}。
成立条件：解足够光滑（$H^{2}$）；网格拟一致。

\field{判据} $n=16/32/64/128$ 三次 $h$ 减半的误差比值落在 $[3.8,4.2]$；
2D 用 $u=x(1-x)y(1-y)$、$f=2y(1-y)+2x(1-x)$，阶 2 但节点\textbf{不}精确。

\field{前置} A5、F2、F4。

\field{常见误解} 以为「1D 节点精确说明我的离散是超收敛的」。它是 1D 的结构性巧合；
照搬到 2D 会得到一个必然失败的判据。
\end{know}

\begin{know}{F7}{CSR 与 SpMV 的复杂度}{未建：linalg/sparse.hpp}
\field{是什么} CSR 用三个数组：\texttt{values}（非零元）、\texttt{colIdx}（列号）、
\texttt{rowPtr}（每行起点，长 $n+1$，\texttt{rowPtr[n]}=\text{nnz}）。

\field{为什么需要} 有限元的 $A$ 每行只有几个非零；稠密存储的 $O(n^{2})$ 与「想解多大的问题」直接冲突。

\field{关键结论} 一次 SpMV 是 $O(\text{nnz})$，而稠密是 $O(n^{2})$；对 1D 三对角 $\text{nnz}=3n-2$，
所以每次迭代 $O(n)$ 而非 $O(n^{2})$。成立条件：\texttt{rowPtr} 单调不减且首末正确；
稀疏没有「结构零」以外的假设。

\field{判据} 与稠密乘逐点对拍（随机 $100\times100$、密度 5%、$\le10^{-14}$）；
三对角手算 \texttt{rowPtr}=$[0,2,5,8,\dots]$；边界形状 $1\times1$、单行、全零行、$\text{nnz}=0$ 不许崩。

\field{前置} F1。

\field{常见误解} 以为「稀疏存储一定更快」。$n$ 小或密度高时，索引开销会反超；
判据是「实际耗时比」而不是「理论复杂度」。
\end{know}

\begin{know}{F8}{Krylov 子空间与共轭方向}{未建：linalg/cg.hpp}
\field{是什么} CG 在 Krylov 子空间 $\mathcal K_k=\mathrm{span}\{b,Ab,\dots,A^{k-1}b\}$ 中找
使 $\norm{e_k}_{A}$ 最小的 $x_k$；方向 $p_k$ 满足 $A$-共轭（$p_i^{H}Ap_j=0$，$i\ne j$）。

\field{为什么需要} 它是「迭代解」的原型：只用一个矩阵-向量乘，不存矩阵，内存 $O(n)$。

\field{关键结论} 三项递推：$r_k=b-Ax_k$、$\alpha_k=\frac{r_k^{H}r_k}{p_k^{H}Ap_k}$、
$x_{k+1}=x_k+\alpha_kp_k$、$p_{k+1}=r_{k+1}+\beta_kp_k$。共轭方向保证「每个方向只走一次」，
故理论上 $n$ 步精确收敛。成立条件：$A$ 对称正定（否则 $\norm{\cdot}_{A}$ 不是范数）。

\field{判据} 与直接解对拍，残差 $\le10^{-10}$；迭代数只记录不设阈值（1D 无预条件约 $O(n)$ 量级）；
把 $\alpha$ 写成反号必须立刻发散——这是注入对照。

\field{前置} F7、L4、N8。

\field{常见误解} 以为「共轭梯度就是最速下降的改进版」。最速下降每步正交于残差、会「锯齿」；
CG 每步 $A$-正交于所有历史方向，所以能在一个 $n$ 维空间里走 $n$ 步精确到达。
\end{know}

\begin{know}{F9}{CG 的收敛率与「为什么必须 SPD」}{未建：linalg/cg.hpp}
\field{是什么} 误差界：
\[
\norm{e_k}_{A}\le2\left(\frac{\sqrt\kappa-1}{\sqrt\kappa+1}\right)^{k}\norm{e_0}_{A},
\]
$\kappa=\lambda_{\max}/\lambda_{\min}$。

\field{为什么需要} 它决定「迭代多少次」，也解释「为什么预处理是必要的」（把 $\kappa$ 压小）。

\field{关键结论} $\kappa$ 大时收敛慢（$\kappa=10^{4}$ 需要约 $100$ 次以上量级）；
非 SPD 时 $p^{H}Ap$ 可能为 0 或负，$\alpha$ 无意义 → 残差停滞或上升。
成立条件：$A$ 对称正定；精确算术；浮点上收敛会因舍入而偏离界。

\field{判据} 非 SPD 的 \texttt{apply}（例如未加边界的纯拉普拉斯）必须表现为残差停滞/上升并被
最大迭代数截断，\textbf{不许死循环}。

\field{前置} F8、A4、N8。

\field{常见误解} 以为「CG 对任何矩阵都能用，只是慢一点」。对非对称矩阵它连「下降方向」都无法保证，
会直接崩；这就是为什么对称性（F5）是硬约束。
\end{know}

\section{张量代数与微分形式}

\subsection*{串讲六：弱形式与 $d^{2}=0$ 是同一条边界的两种写法}

F3 的分部积分和 M3 的 $d^{2}=0$ 看起来一个属于有限元、一个属于微分几何，其实是同一件事的两端：

\begin{center}
\begin{tabular}{L{2.4cm}L{5.0cm}L{6.4cm}}
\toprule
 & 有限元（F3） & 微分形式（M3、M6） \\
\midrule
核心恒等式 & $\int_\Omega\nabla u\cdot\nabla v=-\int_\Omega(\Delta u)v+\int_{\partial\Omega}\partial_nu\,v$ & $\int_M d\omega=\int_{\partial M}\omega$ \\
「挪一次导数」 & 分部积分 & 外微分 + Stokes \\
边界项怎么来 & 被积函数求导产生全导数项 & 定向边界上的低一阶形式 \\
二阶导消失 & 弱形式只含一阶 & $d^{2}=0$：再微一次必为 0 \\
\bottomrule
\end{tabular}
\end{center}

所以「为什么有限元必须用弱形式」与「为什么 $d^{2}=0$」本质是同一个机制：\textbf{求导会把信息推到边界上}。
把这句话记住，两块内容就只需要记一份。

\subsection*{串讲七：三种「连续性」，别混}

\begin{center}
\begin{tabular}{L{2.6cm}L{3.6cm}L{3.4cm}L{3.0cm}}
\toprule
说法 & 属于谁 & 判据 & 违反后果 \\
\midrule
张量连续性（连续布局） & T3 & 步长与形状推出的步长逐维相等 & reshape 必须抛错（T5） \\
函数光滑性 $C^{p}$ & H11 & 第 $p$ 阶导数连续 & 谱系数衰减只有 $O(k^{-p-1})$ \\
收敛阶 $O(h^{p})$ & A5 & $h$ 减半误差降 $2^{p}$ & 渐近区外比值乱、判据假红 \\
\bottomrule
\end{tabular}
\end{center}

三者都叫「连续/光滑」，但一个讲内存、一个讲函数、一个讲误差率。

\begin{know}{T1}{形状与行主序步长}{框架：core/tensor.hpp}
\field{是什么} 秩 $r$、形状 $\text{shape}[0..r-1]$；步长从最后一维往前算：
$\text{strides}[r-1]=1$，$\text{strides}[d]=\text{strides}[d+1]\cdot\text{shape}[d+1]$，
即 $\text{strides}[d]=\prod_{k>d}\text{shape}[k]$。

\field{为什么需要} 步长是「多维下标 → 一维偏移」的全部信息；视图、切片、换轴都只改步长不动数据。

\field{关键结论} 必须\textbf{从后往前}算（行主序）。成立条件：$r\ge1$；$\text{shape}[d]\ge0$。

\field{判据} $\{2,3,4\}\Rightarrow\text{strides}=\{12,4,1\}$。

\field{前置} L1。

\field{常见误解} 以为「步长就是形状的累乘」。只在行主序且连续时才对；
换轴之后步长与形状不再有那个关系（T4），这正是「视图」能零拷贝的原因。
\end{know}

\begin{know}{T2}{下标到偏移的线性映射}{框架：core/tensor.hpp}
\field{是什么} $\text{offset}=\sum_{d=0}^{r-1}\text{idxs}[d]\cdot\text{strides}[d]$。

\field{为什么需要} 它是「多维数组」的全部实现原理：所有访问都是这一行点积。

\field{关键结论} 这是\textbf{线性映射}（对下标线性），所以换轴=换求和顺序、切片=给某一项加常数。
成立条件：$0\le\text{idxs}[d]<\text{shape}[d]$；下标个数必须等于秩。库只在 \texttt{at()} 里查越界，
\texttt{operator()} 不逐维查——个数写错会越界读或丢下标（静默错）。

\field{判据} shape $\{2,3,4\}$ 下 $(1,2,3)=1\cdot12+2\cdot4+3=23$。

\field{前置} T1。

\field{常见误解} 以为「列主序只是把两个下标换一下」。步长整组都变（$\text{strides}=\{1,2,6\}$），
所有切片与换轴的语义随之改变；两套约定不能局部混用。
\end{know}

\begin{know}{T3}{连续性的正确定义}{框架：core/tensor.hpp}
\field{是什么} $\text{IsContiguous}\iff\text{strides}=\text{ComputeStrides(shape)}$（逐维相等）。

\field{为什么需要} 它是 reshape 能否零拷贝的唯一判据；也是性能上「能不能走顺序访问」的判据。

\field{关键结论} 连续性是\textbf{布局性质}，与指针是否相邻无关——换轴后某些情形下标恰好仍然连续，
但一般不是；用「肉眼看着挨着」判断必错。成立条件：秩相同。

\field{判据} 直接比两组步长（不抽样、不看内容）。

\field{前置} T1。

\field{常见误解} 以为「转置之后还是连续」。转置交换了步长（$\{12,4,1\}\to\{1,4,12\}$），
绝大多数情况不再连续；numpy 也是这样。
\end{know}

\begin{know}{T4}{切片与换轴：零拷贝视图}{框架：core/tensor.hpp}
\field{是什么} 沿轴 $a$ 取 $[b,e)$：指针加 $b\cdot\text{strides}[a]$、$\text{shape}'[a]=e-b$、步长不变。
换轴（置换）：$\text{shape}'[d]=\text{shape}[\text{axes}[d]]$、$\text{strides}'[d]=\text{strides}[\text{axes}[d]]$、指针不变。

\field{为什么需要} 大张量的切片必须是 $O(1)$；复制会让「取一个通道」这种操作退化成整块搬运。

\field{关键结论} 两者都不搬字节，所以都\textbf{与宿主共享内存}——视图生命周期不得长于宿主（契约，不做运行期检查）。
成立条件：切片区间合法；置换须是 $0..r-1$ 的排列。

\field{判据} 切片/换轴后元素总数与预期一致；修改视图要能反映到宿主（别名语义）。

\field{前置} T1、T3。

\field{常见误解} 以为「换轴会把数据搬成新布局」。不搬；搬的话 reshape 就不需要「非连续必须抛错」这条规则了。
\end{know}

\begin{know}{T5}{reshape 的合法条件与拷贝顺序}{框架：core/tensor.hpp}
\field{是什么} 合法的充要条件：(1) $\text{IsContiguous}()$；(2) $\prod\text{newShape}=\text{Size}()$。

\field{为什么需要} 非连续张量的「内存顺序」与新形状的行主序不一致，直接重解释会得到\textbf{形状对、数值错}的结果。

\field{关键结论} 深拷贝必须按「里程表」从\textbf{最后一维}进位递增下标，
\[\text{out}[n]=\text{ptr}\big[\text{Offset}(\text{idx},\text{srcStrides})\big];\]
从第 0 维进位会得到转置结果。
成立条件：连续性 + 元素数相等。非连续时库\textbf{一律抛错}（numpy 会静默复制，语义陷阱）。

\field{判据} 拿 $a^{\mathsf T}$ 的拷贝与「转置后的拷贝」对比，两者必须不同（否则进位方向错了）。

\field{前置} T1、T2、T3。

\field{常见误解} 以为「reshape 就是改形状标签」。只有连续张量才可以只改标签；否则必须复制，
而复制顺序错了就是个静默的转置 bug。
\end{know}

\begin{know}{T6}{广播与形状推导}{框架：core/tensor.hpp（语义参照 autodiff）}
\field{是什么} 逐元素运算要求形状相同或可广播（例如 $1\times1$ 的标量与 $m\times n$ 张量）——
库的选择是：\textbf{只在明确的 $1\times1$ 情形放行}，其余同形状要求，否则抛错。

\field{为什么需要} 广播能省内存，但也让形状错误变得难以发现；自动微分的雅可比形状尤其依赖这条约定。

\field{关键结论} 广播的数学后果：$f$ 为标量函数作用在 $1\times1$ 上时，$J=\varphi'(a)\cdot\mathbf 1$（全 1 矩阵）；
形状错误必须\textbf{逐维}检查（只比一个维度会退化成恒真——源码记录过这个 bug 导致越界读）。
成立条件：被广播维度确实是 1，或形状完全一致。

\field{判据} 形状校验逐维；广播情形的 $J$ 与解析值对拍。

\field{前置} T1、D3。

\field{常见误解} 以为「广播是免费的」。它放大的是\textbf{读写次数}（同一份数据被重复读），
在 cache 带宽受限时是实打实的代价。
\end{know}

\begin{know}{T7}{元素总数与空乘积}{框架：core/tensor.hpp}
\field{是什么} $\text{Size}()=\prod_d\text{shape}[d]$；秩 0（标量）时\textbf{空乘积为 1}。

\field{为什么需要} 标量必须有一个元素而不是零个；否则「$1\times1$ 矩阵」与「标量」的运算语义会分裂。

\field{关键结论} 任一维为 0 $\Rightarrow\text{Size}=0$（合法空张量，不是错误）；含 0 维时\textbf{秩仍照算}
（$\{2,0,3\}$ 的秩是 3、元素数 0）。成立条件：形状非负。

\field{判据} 空张量与标量分别用例；`0×0` 的写盘/读回（对应 NOW 的 npy 契约）也要能过。

\field{前置} T1。

\field{常见误解} 以为「元素数为 0 就是秩为 0」。秩是维数、元素数是乘积；$\{0\}$ 的秩是 1、元素数是 0。
\end{know}

\begin{know}{M1}{$k$-形式与分量存放}{未建：math/forms.hpp}
\field{是什么} $k$-形式是反对称的 $k$ 重线性映射；$\R^{n}$ 上分量个数为 $\binom nk$，
用递增下标 $i_1<i_2<\dots<i_k$ 作字典序索引，其余分量由反对称性决定符号。

\field{为什么需要} 它是「有向面积/体积」的代数化；梯度、旋度、散度都是它在 $n=3$ 的不同阶表现（M7）。

\field{关键结论} 反对称性 $\omega(\dots,v_i,\dots,v_j,\dots)=-\omega(\dots,v_j,\dots,v_i,\dots)$ 决定一切；
$k>n$ 时不存在非零 $k$-形式。成立条件：单个坐标卡（不做流形与图册）。

\field{判据} $d^{2}=0$ 与 Leibniz 律在二次多项式系数上逐分量 $\le10^{-14}$（代数恒等式，不是收敛）。

\field{前置} L7、T1。

\field{常见误解} 以为「$k$-形式就是反对称张量」。还要有\textbf{拉回}与\textbf{外微分}这两件配套运算，
否则它只是个带符号的表。
\end{know}

\begin{know}{M2}{外积（楔积）}{未建：math/forms.hpp}
\field{是什么} $(\omega\wedge\eta)(v_1,\dots,v_{p+q})=\sum_\sigma\mathrm{sgn}(\sigma)\,
\omega(v_{\sigma(1)},\dots)\eta(v_{\sigma(p+1)},\dots)/\text{重复计数}$；
分量层面就是「按组合数展开并带符号」。

\field{为什么需要} 它把低阶形式建成高阶，且天然反对称——行列式、体积形式 $dx\wedge dy\wedge dz$ 都是它。

\field{关键结论} 反交换：$\omega\wedge\eta=(-1)^{pq}\eta\wedge\omega$（$p,q$ 为阶数）；
$\omega\wedge\omega=0$ 当 $p$ 为奇数。成立条件：$p+q\le n$（超过就表示层存不下，抛错）。

\field{判据} 与显式叉积/行列式对拍（$n=3$ 的全部组合）；$k+l>n$ 必须抛错而不是截断。

\field{前置} M1。

\field{常见误解} 以为「$\omega\wedge\eta=\eta\wedge\omega$ 差个符号而已」。差的符号是 $(-1)^{pq}$，
两个奇形式交换\textbf{不变号}（$pq$ 为奇 $\times$ 奇 = 奇？注意 $p=q=1$ 时 $pq=1$，故反号）。
算清楚指数，别背口诀。
\end{know}

\begin{know}{M3}{外微分与 $d^{2}=0$}{未建：math/forms.hpp}
\field{是什么} $d$ 把 $k$-形式映到 $k+1$ 形式，分量上是对系数取偏导再反对称组合；
库的实现里\textbf{偏导由调用方给}（$C(n,k)\times n$ 矩阵），只做反对称组装。

\field{为什么需要} 它是「梯度、旋度、散度」的统一算子，也是 $d^{2}=0$ 这条「边界没有边界」的代数表述。

\field{关键结论} 两条恒等式：$d^{2}=0$；Leibniz 律 $d(\omega\wedge\eta)=d\omega\wedge\eta+(-1)^{p}\omega\wedge d\eta$
（$p=\deg\omega$，符号来自反交换）。成立条件：系数 $C^{2}$（混合偏导可交换）。

\field{判据} 在二次多项式上 $d^{2}=0$ 逐分量 $\le10^{-14}$；Leibniz 律同样；
「$d$ 的符号约定」用 $n=3$ 字典（M7）反查，不与叉积冲突。

\field{前置} M1、M2。

\field{常见误解} 以为「$d^{2}=0$ 只是符号巧合」。它的实质是「混合偏导可交换 + 反对称化」，
所以凡是对称二阶导可交换的场合都成立；在离散化里它会以「边界算子复合为零」的形式重现（串讲六）。
\end{know}

\begin{know}{M4}{拉回}{未建：math/forms.hpp}
\field{是什么} 光滑映射 $F:M\to N$ 把 $N$ 上的 $k$-形式拉回成 $M$ 上的 $k$-形式：
$(F^{*}\omega)_p(v_1,\dots,v_k)=\omega_{F(p)}(dF_pv_1,\dots,dF_pv_k)$；分量上是 Cauchy--Binet 子式。

\field{为什么需要} 换元积分、坐标变换下「同一个量怎么写」全靠它；也是「$F^{*}$ 与 $d$ 可交换」这条结构的载体。

\field{关键结论} $F^{*}(d\omega)=d(F^{*}\omega)$（与 $d$ 交换）；$F^{*}(dx\wedge dy)=\det(J)\,dx\wedge dy$（2 维）。
成立条件：$F$ 光滑；定向一致（否则符号翻转）。

\field{判据} 线性 $F$ 与二次 $F$ 各一例；$\det(J)$ 与显式行列式对拍。

\field{前置} M2、M3、L9。

\field{常见误解} 以为「拉回就是换元」。换元是拉回在积分里的\textbf{一个应用}；
拉回本身定义在形式上，与积分无关，所以它才能与 $d$ 交换。
\end{know}

\begin{know}{M5}{Hodge star}{未建：math/forms.hpp}
\field{是什么} 在 $\R^{n}$ 的欧氏度量与标准定向下，$\star$ 把 $k$-形式映到 $(n-k)$-形式，
由 $\omega\wedge\star\eta=\inner{\omega}{\eta}\,\mathrm{vol}$ 唯一确定。

\field{为什么需要} 它把「内积」翻译成「楔积 + 体积」，从而把散度写成 $d\star$，把拉普拉斯写成 $\star d\star d$。

\field{关键结论} 对合性 $\star\star\omega=(-1)^{k(n-k)}\omega$；成立条件：只做欧氏度量（非欧需度量矩阵）。

\field{判据} 对合性逐分量 $\le10^{-14}$；$n=3$ 时 $\star$ 与叉积/散度的对应（M7）。

\field{前置} M2、L4。

\field{常见误解} 以为「$\star$ 就是补下标」。补下标只是一半，还要按体积元与定向确定符号；
搞错符号会让散度差一个负号，在 PDE 里表现为「椭圆算子变双曲」。
\end{know}

\begin{know}{M6}{Stokes 三定理的统一形式}{未建：math/forms.hpp}
\field{是什么} $\int_M d\omega=\int_{\partial M}\omega$。取 $n=3$ 的不同阶即得 Green、散度、Stokes 三条经典定理。

\field{为什么需要} 它把「边界上的积分」与「内部微分」等价，是有限元弱形式（F3）的几何版本。

\field{关键结论} 三条退化：0-形式给梯度定理、1-形式给 Stokes（环量=旋度通量）、
2-形式给散度定理（通量=散度积分）。成立条件：$M$ 有定向、边界分段光滑、$\omega$ 在闭包上 $C^{1}$。

\field{判据} 三次多项式 + 张量积 Simpson（对三次精确），两边差 $\le10^{-13}$。

\field{前置} M3、M4。

\field{常见误解} 以为「这是三个定理」。它是一个定理的三次取值；记住 $\int_{\partial}=\int d$ 即可，
三条公式都能现推。
\end{know}

\begin{know}{M7}{$n=3$ 字典：grad / curl / div}{未建：math/forms.hpp}
\field{是什么} 0-形式 $\leftrightarrow$ 梯度、1-形式 $\leftrightarrow$ 旋度、2-形式 $\leftrightarrow$ 散度；
即 $d$ 在三个阶上分别实现这三个算子（配合 M5 的 $\star$）。

\field{为什么需要} 它是「微分形式不是抽象玩具」的证据：向量分析的三件套就是同一个 $d$ 的不同阶表现。

\field{关键结论} 两条熟知的恒等式在此变成一句话：$\mathrm{curl}\,\mathrm{grad}=0$ 与
$\mathrm{div}\,\mathrm{curl}=0$ 都是 $d^{2}=0$。成立条件：欧氏度量 + 标准定向。

\field{判据} 与显式叉积/散度逐分量对拍；$d^{2}=0$ 的三阶版本各自验证。

\field{前置} M3、M5。

\field{常见误解} 以为「$\mathrm{curl}\,\mathrm{grad}=0$ 是算出来的巧合」。它是拓扑事实
（$d^{2}=0$），与坐标系无关；换成曲线坐标依然成立，而手算展开会累死。
\end{know}

\section{两条轴：离散化与容差}

\noindent\textbf{入口}：库里 97 条数学都活在两条轴上——一条是\textbf{连续 $\to$ 离散}（定理变成有限步算法），
一条是\textbf{精确 $\to$ 容差}（算法落到浮点）。说「这个算法对」之前，必须先说清它在这两条轴上各掉了多少。
本册的每条判据，都是这两条轴上的一个坐标。

\subsection{轴一：离散——定理怎么变成有限步}

连续对象换成有限样本、网格或基函数之后，误差按步长走：

\[
\norm{u-u_h}\le C\,h^{p},
\]

$h$ 是网格尺寸，$p$ 是方法阶。这条不等式是全册最重要的形状：\textbf{它把「对不对」换成了「按什么速率收敛」，后者才是可测的。}

\begin{center}
\begin{tabular}{L{3.4cm}L{5.2cm}L{6.2cm}}
\toprule
离散手段 & 判据 & 失效点 \\
\midrule
网格 + 有限元 & $h$ 减半，误差降 $2^{p}$，$p=$ 基函数多项式次数 & $h$ 未进渐近区时测出的「阶」不是阶 \\
采样 & 采样率 $>2f_{\max}$；折叠频率 $f \bmod f_s$ & 欠采样后混叠不可逆，无法从样本恢复 \\
求积 & $n$ 点 Gauss 对 $2n-1$ 次多项式精确 & 被积函数有间断，阶掉到 1 \\
截断（DFT / 级数） & 残差按 $\abs{\lambda_N}$ 或系数衰减 & 系数的衰减由光滑度决定，不由谱决定 \\
\bottomrule
\end{tabular}
\end{center}

\noindent\textbf{一条经验}：阶的测量必须用\textbf{三个点}（$h$、$h/2$、$h/4$），比值落在 $[2^{p}-\delta,\ 2^{p}+\delta]$ 才算测到；
单点比值乱跳只说明还没进渐近区。

\subsection{轴二：容差——「相等」必须换成「$\le \eps$」}

浮点没有相等，只有相对与绝对容差。所以本册的判据一律写成不等式，且必须带三要素：
\textbf{量、阈值、阈值的量级依据}。缺一条就不写。

\begin{center}
\begin{tabular}{L{4.6cm}L{9.8cm}}
\toprule
判据形状 & 用法与陷阱 \\
\midrule
$\norm{x-y}\le \eps_{\text{mach}}\norm{x}$ & 只对同量级有效；量级差三个数量级时相对误差已经不适用 \\
$\norm{x-y}\le a\,\eps_{\text{mach}}+r\,\norm{y}$ & 绝对 + 相对双阈值，库里比较函数统一用这个形状 \\
$\norm{A^{H}(b-Ax)}\le tol\cdot\norm{A^{H}b}$ & 最小二乘判\textbf{后向}残差；前向误差还要除以条件数 \\
$\kappa(A)\,\eps_{\text{mach}}\ll 1$ & 条件数放大：算法再好也救不了病态问题 \\
累积 $\sqrt{n}\,\eps$ 或 $n\,\eps$ & 随机游走型与最坏情形型，二者差一个 $\sqrt n$ \\
\bottomrule
\end{tabular}
\end{center}

\noindent\textbf{两条曲线}：$h\to0$ 时离散误差按 $h^{p}$ 降，但舍入被 $\kappa=O(h^{-2})$ 放大，
总误差在双对数图上出现\textbf{谷底}。\textbf{收敛阶测试的三个点必须都在谷底左侧}，否则测出来的不是 $p$。

\subsection{两条轴的交叉：为什么「数学正确」不等于「程序对」}

\begin{center}
\begin{tabular}{L{2.8cm}L{3.0cm}L{4.6cm}L{3.8cm}}
\toprule
定理 & 离散实现 & 浮点判据 & 失效点 \\
\midrule
正交归一 & MGS & $\norm{Q^{H}Q-I}\le 8n\eps$ & 不做重正交时正交性随 $n$ 退化 \\
可逆 & LU + 部分主元 & 残差小 & 残差小但解错：条件数大 \\
精确求积 & Simpson & 三次多项式精确 & 间断处阶降到 1 \\
谱分解 & QR 迭代 & 残差 $\le tol$ & 重根时特征向量不稳 \\
\bottomrule
\end{tabular}
\end{center}

\noindent\textbf{这张表的读法}：左边是数学，中间是实现，右边是判据，最后一列是判据失效的地方。
\textbf{四条缺一条就不算验证过}——只给左边的，是教材；只给中间的，是代码；四列齐全的，才是本册的条目。

\subsection{与「入口法」的关系（只借用两格）}

根目录 \texttt{LEARNING.md} 的六段在本册只用两格：
\textbf{入口}（这条数学在代码里长成什么形状）与\textbf{失效点}（状态词与数值极限）。
另外四格里的「生成关系」在本册有确定含义：\textbf{连续 $\to$ 离散 $\to$ 浮点}，每换一步就掉一次精度，
而每一掉都在上面两张表里有一行对应。这就是本册不写推导、只写判据的理由。

\section{覆盖度自查与缺口}\label{sec:audit}

\subsection{文件到条目的映射}

下表是「每个文件贡献了哪几条」。\textbf{没有任何一个文件是零贡献}——
最少的 \texttt{core/common.hpp} 与 \texttt{core/scalar\_traits.hpp} 也各有归属
（前者给 $\eps$ 的实际取值与 $\pi$ 精度口径，后者自述「没有数学公式」，故并入 A1 的口径说明）。

\begin{longtable}{L{4.8cm}L{1.2cm}L{2.7cm}L{5.2cm}}
\toprule
文件 & 状态 & 贡献条目 & 归不进数学的内容 \\
\midrule
\endhead
\texttt{core/scalar.hpp} & 已实现 & A1、A2、L4 & concept 约束写法、静态函数替代 constexpr 的工程理由 \\
\texttt{core/complex.hpp} & 已实现 & C1--C6、A6 & 不用 \texttt{std::complex} 的三条理由 \\
\texttt{core/matrix.hpp} & 已实现 & L1--L3、L10 & 哑容器设计、i-k-j 循环的 cache 友好性 \\
\texttt{core/tensor.hpp} & 框架 & T1--T7 & 视图生命周期契约、不做并发/表达式模板 \\
\texttt{core/common.hpp} & 部分注释掉 & A1（$\eps$ 取值）、A2（$\pi$ 口径） & 三个 enum 的结构标签 \\
\texttt{core/scalar\_traits.hpp} & 已实现 & （并入 A1 说明） & 一行前向声明用于打断循环包含 \\
\texttt{linalg/solve.hpp} & 已实现 & N1、N2、A2、A3 & 行交换簿记、\texttt{piv} 存列号 \\
\texttt{linalg/inverse.hpp} & 已实现 & N9、N8、L3 & 五条路径不复用的工程理由 \\
\texttt{linalg/qr.hpp} & 已实现 & N3、N4、N5、N6、L10 & 三条路径对拍策略、\texttt{Adjoint} 踩坑记录 \\
\texttt{linalg/cholesky.hpp} & 已实现 & N7、N8 & 异常类型的历史遗留 \\
\texttt{linalg/svd.hpp} & 已实现 & L7、N12、N13、A4 & 排序须与 $V$ 列同步置换 \\
\texttt{linalg/eigen.hpp} & 已实现 & L6、L8、N11 & 特征值返回顺序不保证 \\
\texttt{linalg/expm.hpp} & 已实现 & N10、C3 & 求逆 vs 求解的 flops 纪律 \\
\texttt{linalg/operator.hpp} & 已实现 & L4、L5 & 暂驻/搬家理由、禁用 \texttt{+=} \\
\texttt{linalg/so3.hpp} & 已实现 & G1--G4 & 三层否决的分层理由、不查反对称 \\
\texttt{quat/quaternion.hpp} & 已实现 & G5--G10 & 不提供标量除法的二义性 \\
\texttt{autodiff/dual.hpp} & 已实现 & D1、D2 & ADL/友元断链、无复合赋值 \\
\texttt{autodiff/jacobian.hpp} & 已实现 & D3、D4 & 重载偏序、形状检查的旧 bug \\
\texttt{autodiff/grad\_graph.hpp} & 已实现 & D5、L10 & tape 的拓扑排序实现 \\
\texttt{math/ode.hpp} & 骨架 & D6--D8 & \texttt{std::function} 别名、异常分类 \\
\texttt{math/fourier.hpp} & 框架 & H1--H8 & 长度契约分档、不自动分派 \\
\texttt{math/series.hpp} & 框架 & H9--H11 & 与 fourier 分文件的理由 \\
\texttt{math/special.hpp} & 框架 & P8--P10、A6 & 默认参数、不用 Dual/Complex 的取舍 \\
\texttt{math/prob.hpp} & 框架 & P1--P7 & seed 显式化、逐位可复现口径 \\
\texttt{NOW: grid.hpp} & 未建 & F1 & 手算表对拍的验收方式 \\
\texttt{NOW: sparse.hpp} & 未建 & F7 & 「生产路径不许 ToDense」的纪律 \\
\texttt{NOW: assemble.hpp} & 未建 & F2--F6 & 非多项式 $f$ 需要求积器（不做） \\
\texttt{NOW: cg.hpp} & 未建 & F8、F9 & \texttt{std::invocable} 零开销约束 \\
\texttt{NOW: forms.hpp} & 未建 & M1--M7 & 「定义用性质、公式作推论」的书写纪律 \\
\bottomrule
\end{longtable}

\subsection{97 条的构成}

\begin{center}
\begin{tabular}{L{3.0cm}L{2.0cm}L{9.0cm}}
\toprule
类别 & 条数 & 含义 \\
\midrule
已实现 & 49 & 有真实算法，判据可跑（A1--A6、L1--L8、L10、N1--N13、C1--C6、D1--D5、G1--G10）\\
框架 & 28 & 公式可信、函数体全 throw（H1--H11、P1--P10、T1--T7）\\
骨架 & 3 & 源码自述骨架（D6--D8）\\
未建 & 16 & 只在 \texttt{NOW.md} 规划（F1--F9、M1--M7）\\
缺位 & 1 & 库里没有位置也没规划（L9 行列式）\\
\midrule
合计 & 97 & \\
\bottomrule
\end{tabular}
\end{center}

\textbf{这个分布本身就是一条结论}：这张清单里 31 条（框架+骨架）的公式是「设计意图」而非「已验证事实」，
16 条连文件都还没有。它们能用来复习，\textbf{不能}用来论证「库实现了这个数学」。

\subsection{缺口：库用到、但数学上没成体系}

以下六处是真正值得补的空白。它们的共同点：\textbf{不是缺一条公式，而是缺一段推导}。

\begin{enumerate}[leftmargin=1.6em]
  \item \textbf{$8n\eps$ 背后的舍入分析}（A2）。源码只说「系数 8 是余量、非推导」。
        完整的推导需要：$n$ 步运算的误差累积上界、以及「阈值反推」四步（$R=e_{\text{tol}}/e_{\text{worst}}\in[10^{1},10^{6}]$）。
        这一段在 \texttt{03-test-standard.md} 里只有工程口径，没有数学证明。

  \item \textbf{条件数 $\kappa(A)$ 的定义与增长}（A4、N12）。\texttt{svd.hpp} 之外，
        其余文件的「三、公式与符号表」都\textbf{没有}出现 $\kappa$。这是全库最明显的一处「用了但没写」。

  \item \textbf{SVD 走 $A^{H}A$ 的精度损失链条}（N12）。$\kappa(A^{H}A)=\kappa(A)^{2}$ 只是一句话；
        「$\sigma$ 的相对误差为何是 $\sqrt\eps$ 量级」还差一步：$\lambda$ 的误差 $\eps\kappa^{2}$ 经开方变成
        $\sqrt{\eps}\kappa$。Golub--Kahan 双对角化能避开，但未实现。

  \item \textbf{Gibbs 过冲常数 $\frac2\pi\int_0^\pi\frac{\sin t}{t}dt$ 的来源}（H10）。
        数值 $1.08949\ldots$ 有了，但「为什么是 sinc 积分、为什么与 $K$ 无关」的推导没有。

  \item \textbf{QR 与最小二乘的联系}（N6）。\texttt{qr.hpp} 的「三、公式与符号表」里\textbf{完全没有}
        最小二乘；而最小二乘是 QR 最重要的消费场景。这条是本清单里唯一「数学上必需、源码里缺失」的。

  \item \textbf{行列式与体积}（L9）。库没有 \texttt{det} 入口，但拉回（M4）需要 $\det J$，
        第四节的「体积形式」也需要它——目前只能靠 SVD 的恒等式绕过去。
\end{enumerate}

\subsection{另外三处「形式上完整、实质待验证」}

\begin{itemize}[leftmargin=1.4em]
  \item \textbf{框架层的判据全是意图}：H、P、T 三组的阈值（$10^{-13}$、$2$ ulp、$10^{-14}$）都写在注释里，
        没有一条跑过。它们的价值是「给出了验收目标」，不是「已经验收」。
  \item \textbf{骨架层的阶数是承诺}：D7 的「RK4 四阶」、D8 的「能量有界」都还没有实现，
        所以「辛 vs 非辛」的对照目前只是数学陈述。
  \item \textbf{未建层的编号是计划}：F1 的「$4\times3$ 网格有 12 个三角形」这类数字是手算目标，
        库里还没有能跑的代码去核对它。
\end{itemize}

\subsection{最后一句}

这份清单能保证的是：\textbf{每个数学对象都找到了它在库里的位置，且说清了它的成立条件与失败长相}。
它不能保证的是：这些条件的\textbf{数值}已经验证过。二者的差别，就是「读过」与「做过」的差别。


\end{document}
